\documentclass[11pt]{amsart}
\usepackage[margin=1in]{geometry}
\usepackage{amsmath,amssymb,amsthm,mathtools,booktabs,longtable,array,microtype}
\usepackage{hyperref}
\hypersetup{colorlinks=true,linkcolor=blue,citecolor=blue,urlcolor=blue,
 pdftitle={A Lower Exponent of 1.0418235 for Planar Unit Distances},
 pdfauthor={Eric Naslund}}
\newtheorem{theorem}{Theorem}[section]
\newtheorem{proposition}[theorem]{Proposition}
\newtheorem{lemma}[theorem]{Lemma}
\newtheorem{corollary}[theorem]{Corollary}
\theoremstyle{definition}
\newtheorem{definition}[theorem]{Definition}
\theoremstyle{remark}
\newtheorem{remark}[theorem]{Remark}
\DeclareMathOperator{\Gal}{Gal}
\DeclareMathOperator{\Cl}{Cl}
\DeclareMathOperator{\rd}{rd}
\DeclareMathOperator{\Tr}{Tr}
\DeclareMathOperator{\rank}{rank}
\newcommand{\Ftwo}{\mathbb F_2}
\newcommand{\Q}{\mathbb Q}
\newcommand{\R}{\mathbb R}
\newcommand{\C}{\mathbb C}
\newcommand{\Z}{\mathbb Z}
\newcommand{\deltastar}{0.0418235}
\newcommand{\thetastar}{4095/8192}
\newcommand{\ceiling}{0.042161819}
\numberwithin{equation}{section}
\setcounter{tocdepth}{1}
\raggedbottom
\title[A Lower Exponent for Planar Unit Distances]{A Lower Exponent of 1.0418235\\for Planar Unit Distances}
\author[Eric Naslund]{Eric Naslund\textsuperscript{*}}
\email{naslund.math@gmail.com}
% The AMS title-note slot places the disclosure below the author, before the abstract.
\dedicatory{\normalfont\parbox{0.9\textwidth}{\footnotesize
\textsuperscript{*}This paper was entirely written with AI tools. The listed
human author is responsible for having supervised and prompted the AI agents
that wrote this paper.}}
\date{September 27, 2026}
\begin{document}
\begin{abstract}
Let $u(U)$ count the unordered pairs at distance one in a finite planar
set $U$. Following the number-field construction made explicit by Sawin,
we construct sets with $|U_j|\to\infty$ and
\[
 \frac{u(U_j)}{|U_j|^{1.0418235}}\longrightarrow\infty.
\]
We use mixed-signature fixed fields of complex conjugation in a pro-$2$
tower with prescribed nonabelian dyadic completion. Averaging norm-one
unit orbits cancels the relative regulator and capitulation loss.
A complete Hecke-factor calculation in a degree-$16384$ subfield
controls the relative zeta value. Correlated archimedean profiles
and weighted finite-place shells give the final improvement.
The numerical inequalities are certified by exact finite data and
directed interval bounds.
\end{abstract}
\maketitle
\tableofcontents
\clearpage
\section{Introduction}
For a finite set $U\subset\R^2$, let
\[
 u(U)=\#\{\{x,y\}\subset U:|x-y|=1\},\qquad
 u(n)=\max_{|U|=n}u(U),
\]
where the pairs are unordered and the distance is Euclidean.
Erd\H{o}s introduced the unit distance problem in 1946
\cite{Erdos1946}. His lattice construction gives
\[
 u(n)\ge n^{1+c/\log\log n}
\]
for some $c>0$ and arbitrarily large $n$. The upper bound
\[
 u(n)=O(n^{4/3})
\]
is due to Spencer, Szemer\'edi and Trotter
\cite{SpencerSzemerediTrotter1984}, with a later crossing-number proof
due to Sz\'ekely \cite{Szekely1997}.

The recent number-field construction of OpenAI \cite{OpenAI2026},
explained by Alon et al.~\cite{Alon2026}, gives
$u(n)\ge n^{1+\delta}$ for a fixed $\delta>0$ and arbitrarily large $n$.
Sawin \cite{Sawin2026} made this exponent explicit. In our convention
for unordered pairs, his result gives
\[
 u(n)\ge C^{-1}n^{1.014114}
\]
for an absolute constant $C$ and arbitrarily large $n$.
Emmerich \cite{Emmerich2026} obtained a certificate with
$\delta=0.015263\ldots$ by optimizing the integer parameters in Sawin's
criterion. A subsequent certificate of Emmerich and Cordella
\cite{EmmerichCordella2026} reports $\delta=0.031185334443\ldots$.
The author's earlier manuscript \cite{NaslundLegacy2026} records
$\delta=0.0358324$. We compare its construction with the present one below.

Our main result is:
\begin{theorem}\label{thm:main}
There are finite sets $U_j\subset\R^2$, with $|U_j|\to\infty$, such that
\[
 \frac{u(U_j)}{|U_j|^{1.0418235}}\longrightarrow\infty.
\]
In particular, $u(n)\ge n^{1.0418235}$ for arbitrarily large integers $n$.
\end{theorem}

We obtain the improvement by changing the number fields and counting
norm-one unit orbits in mixed signature, then optimizing the local and
archimedean overlap in the same construction. The result concerns an
unbounded sequence of cardinalities. No generalized Riemann hypothesis
or asymptotic class-number equality is used.

\subsection{Sawin's Construction}

The arithmetic source of unit distances is a quadratic extension
$K/F$ with involution $\iota$. If a prime $\mathfrak p$ of $F$ splits as
$\mathfrak P\iota\mathfrak P$, then the $k+1$ ideals
\[
 \mathfrak P^j(\iota\mathfrak P)^{k-j},\qquad 0\le j\le k,
\]
have the same relative norm $\mathfrak p^k$. Choices at different primes
multiply. An ideal-class argument converts some of these choices into
elements of a common relative norm. At an embedding where $\iota$
acts as complex conjugation, their images have a common length.
After rescaling, they give unit displacements in a planar projection
of an ideal lattice.

Sawin already uses fractional ideals, relative class numbers, and
powers of Frobenius elements to obtain small residue degrees
\cite{Sawin2026}. His explicit construction takes $F$ totally real and
$K/F$ CM. A sup-norm window controls the number of vertices and the
overlap under unit displacements, while Louboutin's relative
class-number bound \cite{Louboutin2000} controls the loss in selecting
an ideal class. Prescribing the local behavior requires relations in
a pro-$2$ group. The Golod--Shafarevich criterion
\cite{GolodShafarevich} ensures that sufficiently few such relations
leave an infinite tower. The use of Frobenius cuts to control splitting
belongs to the class-field-tower methods of Hajir, Maire and Ramakrishna
\cite{HajirMaireRamakrishna2021Cutting}.

There are two related arithmetic precedents. Small split primes and
ideal-class selection occur in the work of Ellenberg and Venkatesh
\cite{EllenbergVenkatesh2007} on class-group torsion. Number fields
with controlled discriminants have also been used to construct codes
and high-dimensional sphere packings
\cite{Lenstra1986,LitsynTsfasman1987}. Here the required lattice
vectors must have the same length after projection to the plane,
which is supplied by their common relative norm.

\subsection{The Earlier Bound $\delta=0.0358324$}

The earlier manuscript \cite{NaslundLegacy2026} keeps the totally
real/CM setting but changes several estimates in Sawin's criterion.
For Euclidean balls of radius $R>1/2$ and a unit vector $e_1$, the
overlap has exponential rate
\[
 \frac{1}{2d}\log
 \frac{\operatorname{vol}_{2d}(B_R\cap(B_R-e_1))}
      {\operatorname{vol}_{2d}(B_R)}
 \longrightarrow \frac12\log\left(1-\frac1{4R^2}\right).
\]
This replaces the loss $\log(1-1/R)$ in the earlier window estimate.
A bound for the dual lattice, together with Banaszczyk's transference
theorem \cite{Banaszczyk1993}, replaces the packing estimate for the
number of vertices by a covolume estimate. The narrow ideal-class
calculation also removes an exponential sign loss.

The analytic estimate uses a fixed multiquadratic subfield of every
field in the tower. Its normalized Euler product retains more
information than a comparison with a power of the Riemann zeta
function. This develops the small-prime comparisons in relative
class-number estimates \cite{Louboutin2000,Louboutin2005}.
The local construction allows dyadic ramification and charges some
relations at degree three in the filtration. For its set of the first
38 odd primes, the displayed infinitude certificate is
\[
 P(t)=1-38t+355t^2+114t^3,\qquad
 P(6/115)=-\frac{101}{1520875}<0.
\]
These changes, followed by reoptimization, give the earlier manuscript's
stated exponent $1.0358324$.

Some of these ideas were described in the author's MathOverflow
answer \cite{NaslundMO2026}, which credits spiderduckpig with placing
the second split-prime Frobenius relation in filtration degree at least
three. Sawin's comment there
also records that he knew the fixed-subfield analytic comparison,
without having evaluated it. The answer and the archived manuscript
are distinct records. In particular, the answer's June 9 change log
reports a bookkeeping correction and the figure $\delta>0.0357$.
We use the archived manuscript's $0.0358324$ to identify the earlier
construction, rather than as a hypothesis of the present proof.

The further improvement comes from allowing a different signature
and different local groups. Covolume counting, ball overlap and the
removal of the sign loss are already part of the earlier argument.
The new counting problem is to use a positive-rank group of norm-one
units without losing its regulator.

\subsection{Mixed Signature and Norm-One Units}

Let $K/\Q$ be totally imaginary and Galois, choose actual complex
conjugation $\iota$, and set $F=K^{\langle\iota\rangle}$. Write
\[
 [K:\Q]=2d,\qquad [F:\Q]=d=b+2c,\qquad \theta=c/d,
\]
where $(b,c)$ is the signature of $F$. The chosen complex embedding of
$K$ restricts to a real embedding of $F$ and supplies the planar
projection. The other embeddings
need not be real, since $\iota$ need not be central. In fact,
\[
 \frac bd=\frac1{|\operatorname{class}(\iota)|},
\]
where the conjugacy class is taken in $\Gal(K/\Q)$. Retaining a finite
quotient with many conjugates of this involution gives a field whose
signature is predominantly complex.

In the CM case, a norm-one element has modulus one at every
archimedean coordinate and the norm-one unit group is finite.
Above a complex place of $F$, its two coordinates can instead have
moduli $e^u$ and $e^{-u}$. The norm-one units have rank $c$ and
translate these reciprocal pairs. We average over a fundamental
domain of their full logarithmic lattice. Tiling gives an exact
integral divided by the relative regulator.

For the extensions used here, which are unramified at every finite
place, the class-number formula cancels this regulator together with
the kernel of extension of ideal classes. More precisely, if
$h_{\rm rel}=h_K/h_F$, $\kappa$ is the order of that kernel,
$R_U$ is the norm-one unit covolume, $w_K$ counts roots of unity, and
$\lambda=\rd(K)$, Proposition~\ref{geo:mass} gives
\[
 \frac{w_K}{h_{\rm rel}\kappa R_U}
 =\frac{2^{d-1}\pi^{b+c}}{\lambda^{d/2}L_F(1)},
 \qquad L_F(s)=\frac{\zeta_K(s)}{\zeta_F(s)}.
\]
At $s=1$, the quotient denotes the quotient of residues.
The averaging uses the entire unit lattice, so no independence of
its coordinates or equidistribution of unit orbits is needed.

\subsection{Local Relations and the Euler Calculation}

To construct the fields, we use the filtered images of the local
relation modules. Fox derivatives \cite{Fox1953} retain the
dependencies among relators that a count of their degrees alone
would discard. The filtration is the one in the
Golod--Shafarevich method, with its group-algebra description
\cite{Jennings,Quillen,Ershov2012}. In particular, a relation shared
by a global condition and a local condition is subtracted only after
identifying the same row in their actual filtered images.

This permits a nonabelian decomposition group of order $32$ at $2$,
with inertia order $8$ and residue degree $4$. Its relation calculation
allows residue degree $2$ at both $3$ and $5$. The resulting family has
\[
 \lambda=2^{9/4}\sqrt{15015},\qquad
 \theta\ge\frac{4095}{8192}.
\]
The gain at $3$ and $5$ has costs. Compared with the preceding
five-prime family, the root discriminant increases from
$4\sqrt{15015}$, and the residue degree at $17$ increases from $2$
to $4$. The smaller residue degrees also increase the corresponding
absolute Euler factors. The exponent improves because the additional
norm choices outweigh these losses in the final inequality.

For the analytic calculation we retain two fields
\[
 N\subset M\subset K,\qquad
 [N:\Q]=16384,\qquad [M:\Q]=524288.
\]
The full Euler calculation is made in $N$. Its regular representation
contains $128$ linear, $256$ quadratic, $456$ quartic and $124$ octic
irreducible representations, and we evaluate their Hecke factors
with the required multiplicities. Further Frobenius information from
$M$ is charged only on the remaining primes. This is necessary because
adding savings from different subfields can otherwise count the same
prime more than once.

The evaluation is at $\sigma=12001/12000$. The unconditional
Tsfasman--Vl\u{a}du\c{t} prime inequality \cite{TsfasmanVladut2002}
then gives, for all sufficiently large fields in the tower,
\[
 \frac1d\log L_F(1)<0.042161819.
\]
Thus the fixed-subfield idea from the earlier paper is used with a
nonabelian field and a complete factor calculation. The finite
field identities, conductor data and analytic remainders are proved
in sections \ref{ar:field} and \ref{an:section}.

\subsection{Archimedean Profiles and Finite Shells}

At a complex place of $F$, a norm-one displacement couples two
coordinates through reciprocal scaling. We consequently use a
profile on $\C^2$ that depends jointly on their radii. It is a positive
cubic Bernstein modification of a Student profile. The coordinates
above real places use a Gaussian. Products are taken over places,
while the two coordinates at each complex place remain coupled.

At a split finite place of residue-field size $Q$, a hard valuation
window gives the factor
\[
 (k+1)Q^{-\delta k}.
\]
We replace it by a weighted sum of six valuation shells. The shell
overlap matrix is explicit, so the gain and the mass can be optimized
together. Section \ref{fw:section} gives the exact finite functional.

Both choices enter the same arithmetic and geometric inequality.
For hard finite windows its form is
\[
 \sum_q\frac{\log(k_q+1)-\delta k_qf_q\log q}{e_qf_q}
 -(1/2-\delta)\log\lambda-C+\mathcal A_\delta(\theta)>0,
\]
where $e_q,f_q$ are absolute local indices, $C$ bounds
$d^{-1}\log L_F(1)$, and $\mathcal A_\delta$ is the archimedean
contribution. Proposition~\ref{geo:transfer} gives the precise
expression, and Proposition~\ref{fw:transfer} replaces the hard
windows by shell weights. Changing the local fields changes several
terms at once, which is why the tower and the profiles are optimized
together.

The weights are converted into an ordinary finite set by restricting
the combined energy to a concentration window. Ideal-class selection
uses the resulting overlap weights. A uniform Poisson estimate then
bounds the vertices for the same translate and unit deformation that
supplies the edges. At
\[
 \delta=\frac{83647}{2000000}=0.0418235,
\]
the rational witness in section \ref{sec:certificate} leaves a
certified margin greater than $5.33204359816794\times10^{-6}$
after the concentration allowance.

\subsection{Computation and Formalization}

The proof uses finite linear algebra, exact coefficient recurrences,
prime data and directed interval enclosures. The arithmetic inputs
and the algorithms for checking them accompany the manuscript.
Section \ref{sec:certificate} distinguishes computations rerun by
the supplied verifier from stored finite data. These certificates
concern numerical inequalities and explicit arithmetic, rather than
the statistical accuracy of a floating-point approximation.

There is also a separate Lean development, using Lean~4 and Mathlib
\cite{Lean4_2021,Mathlib2020}. Its selected theorem proves the planar
conclusion from one explicit inequality for the ordinary Dedekind
zeta function of the fixed field $M$. That inequality is stated in
Remark~\ref{an:fixed-field-form}. It has not been proved in Lean.
The formalization is consequently conditional, and is not used as a
substitute for the analytic argument or finite certificate here.

\section{The Arithmetic Tower}\label{tw:tower}
We construct an infinite pro-$2$ extension with prescribed local groups.
Its finite subextensions will supply the fields in the geometric and
analytic arguments. All pro-$2$ subgroups and normal closures are topological. For a pro-$2$
group $H$, write $D_nH$ for its Zassenhaus filtration, characterized in a
finite quotient by $D_nH=\{h:h-1\in I_H^n\}$, where $I_H$ is the
augmentation ideal of $\mathbb F_2[H]$. The associated restricted Lie
algebra has bracket induced by commutators and restricted square induced
by squaring. Our commutator convention is $[x,y]=x^{-1}y^{-1}xy$.
Put
\[
 T=\{3,5,7,11,13\},\qquad \Sigma=T\cup\{2\},\qquad
 E=\mathbb Q(i,\sqrt2,\sqrt3,\sqrt5,\sqrt7,\sqrt{11},\sqrt{13}).
\]
Let $\mathcal G$ be the Galois group of the maximal pro-$2$ extension of
$\mathbb Q$ unramified at finite primes outside $\Sigma$, with no
restriction at infinity. Fix actual complex conjugation $c$.

\subsection{The Complete Global Presentation}
We use the presentation and cohomological framework for pro-$p$
extensions described in \cite{Koch2002,NeukirchSchmidtWingberg2008}.
The odd local relations and the relation at infinity give a complete
presentation of the global group.
\begin{lemma}\label{tw:complete-global-presentation}
The group $\mathcal G$ has a minimal presentation with $7$
generators and $6$ relators.  The relators can be chosen to be $c^2$ and
the five odd tame local relators.  In particular, the image in the global
free group of any dyadic local relation belongs to their closed normal
closure.
\end{lemma}

\begin{proof}
Set $X=\operatorname{Spec}\mathbb Z[1/\prod_{p\in\Sigma}p]$.
Since $2$ is invertible on $X$, the sheaf $\mu_2$ is constant and Kummer
theory gives
\[
 H^1(X,\mathbb F_2)=
 \langle-1,2,q:q\in T\rangle_{\mathbb F_2}.
\]
The displayed squareclasses are independent, so the generator rank is $7$.

To bound the relation rank, take the inverse system of connected finite
Galois $2$-covers $Y\to X$.
Every class in $H^1(Y,\mathbb F_2)$ is a quadratic torsor.  Its normal
closure over $X$ is again a finite Galois $2$-cover: its Galois group embeds
in an extension of $\operatorname{Gal}(Y/X)$ by a product of copies of
$C_2$.  The direct limit of these first cohomology groups is thus zero. The
five-term sequence of the continuous Cartan--Leray spectral sequence
gives an injection
\begin{equation}\label{tw:global-h2-injection}
 H^2(\mathcal G,\mathbb F_2)
 \lhook\joinrel\longrightarrow H^2(X,\mu_2).
\end{equation}
We obtain the continuous sequence from finite Galois covers, using
the spectral sequence for schemes
\cite[Theorem~14.9]{MilneEtale}.

The Kummer sequence and $\operatorname{Pic}(X)=0$ identify the group on the
right with $\operatorname{Br}(X)[2]$.  Since $X$ is a regular integral
arithmetic curve with finite residue fields, the Brauer residue sequence
identifies $\operatorname{Br}(X)$ with the classes in
$\operatorname{Br}(\mathbb Q)$ having zero invariant outside
$\Sigma\cup\{\infty\}$ \cite[Theorem~6.8.3 and
Example~6.9.4]{Poonen}. Global Brauer reciprocity then identifies its
$2$-torsion with the vectors of local $2$-torsion invariants supported
on $\Sigma\cup\{\infty\}$ whose sum is zero.  In particular its dimension
is $6$. The odd-place and infinite-place coordinates are free, and
reciprocity determines the dyadic coordinate
\cite[Chapter~IV, Example~2.14(e)]{MilneCFT}.

To prove surjectivity, we use the quaternion classes obtained by taking
cups of global quadratic characters.
The class $(-1,-1)$ supplies the infinite coordinate.  If $q\equiv3\pmod4$,
then $(-1,q)$ supplies the coordinate at $q$ among the odd places and
has zero infinite coordinate.  If $(2/q)=-1$, the class $(2,q)$ does the
same.  The remaining primes of $T$ are $5,13$, and $2$ is a quadratic
nonresidue at both. These classes span the odd and infinite coordinates. Their only other
possible nonzero invariant is at two, whose value is determined by
reciprocity.

It follows, by naturality of restriction, that
\begin{equation}\label{tw:local-h2-isomorphism}
 H^2(\mathcal G,\mathbb F_2)
 \xrightarrow{\ \sim\ }
 \bigoplus_{q\in T}H^2(\mathcal G_q,\mathbb F_2)
 \oplus H^2(C_2,\mathbb F_2),
\end{equation}
where $\mathcal G_q$ is the maximal pro-$2$ local group.
At odd $q$, this group has the tame presentation
$\langle\tau_q,\phi_q\mid
\phi_q\tau_q\phi_q^{-1}=\tau_q^q\rangle$ and a one-dimensional second
cohomology group.  The real group is $C_2$ and also has one relation.

For a minimal free pro-$2$ presentation
$1\to R\to\mathcal F\to\mathcal G\to1$, transgression is the
natural perfect pairing
\[
 H^2(\mathcal G,\mathbb F_2)
 \simeq\operatorname{Hom}_{\mathrm{cts}}
       (R/R^2[R,\mathcal F],\mathbb F_2).
\]
This follows from the five-term sequence since $\mathcal F$ is free
\cite[(3.3)--(3.4)]{Quadrelli2021}.
Lift the local generators to $\mathcal F$.  Naturality of transgression
identifies the dual of \eqref{tw:local-h2-isomorphism} with the map taking
each local defining relator to its class in $R/R^2[R,\mathcal F]$.
The odd tame relators and $c^2$ consequently form a basis of this quotient.
The lift of $c$ may be chosen as a free generator, because its image is
nonzero modulo the Frattini subgroup, detected by $\mathbb Q(i)$.

A basis of the relation quotient also gives closed normal generation.
Suppose the indicated relators normally generated $R'\ne R$.
Some finite $2$-quotient of $\mathcal F/R'$ would then have nontrivial image
$N$ of $R/R'$.  The finite nonzero module $N/\Phi(N)$ has nonzero
coinvariants under this finite $2$-group: its group algebra has nilpotent
augmentation ideal.  Hence $N/N^2[N,\mathcal F/R']\ne0$, contradicting
that the indicated relators span $R/R^2[R,\mathcal F]$.
Finally, a dyadic local relation maps to the identity in
$\mathcal G$, so its lift belongs to $R$.
\end{proof}

\subsection{The Prescribed Dyadic Field}
We choose the following extension at two.
\begin{lemma}\label{tw:new-local-two}
There is a Galois extension $L_2/\mathbb Q_2$ with decomposition group
\[
 D=\langle x,y,z\mid x^2,y^2,z^4,[y,z]^2,
                  [x,y],[x,z],[[y,z],z]\rangle
\]
of order $32$, inertia $C_2^3$, residue degree $4$, and normalized
different exponent $9/4$. Its augmentation Hilbert polynomial is
\[
 P_D(t)=(1+t)^3(1+t^2)^2,
\]
and its socle has degree seven.
\end{lemma}
\begin{proof}
Let $U_4/\mathbb Q_2$ be the unramified quartic extension and set
\[
 A=\mathbb Q_2(i,\sqrt5,\sqrt{1+2i}),\qquad
 L_2=A(\sqrt2)U_4.
\]
The norm identity $(1+2i)(1-2i)=5$ gives the dihedral extension
$A/\mathbb Q_2$ of order eight. Its maximal abelian subextension is
$\mathbb Q_2(i,\sqrt5)$ and its unramified degree is two. Thus
$A(\sqrt2)$ has group $D_8\times C_2$, degree sixteen, and unramified
degree two: its elementary abelian abelianization cannot have a cyclic
quotient of order four. Adjoining $U_4$ doubles the degree and the
unramified degree. Its inertia has order eight and is elementary abelian.

Write $w=[y,z]$. The inertia and Frobenius generators can be chosen with
the indicated elementary reciprocity classes $5,-1,-2$. In particular
$z$ is a lift of a uniformizer with the specified unit component.
The elements $x,w$ are central involutions, $y$ is an involution, and
$z$ has order four with $z^2$ independent of $w$. These facts give the
displayed presentation. Conversely, that presentation makes $w$ central:
the identity $w^y=w^{-1}$ and the imposed $w^2=1$ give commutation with
$y$, while commutation with $z$ is imposed explicitly. Its normal forms
$x^a y^b z^c w^d$, with $a,b,d\in\mathbb F_2$ and
$c\in\mathbb Z/4$, have multiplication
\[
 (a,b,c,d)(A,B,C,D)=(a+A,b+B,c+C,d+D+cB).
\]
These normal forms give precisely thirty-two elements. The two
restricted layers have dimensions three and two, which gives the
asserted Hilbert polynomial.

For the different, put $k=\mathbb Q_2(i)$, $\pi=1+i$, and
$\alpha=1+2i=1+\pi^2$. Here $v_k(\pi)=1$ and $v_k(2)=2$.
For $b=1+\pi$, the element
\[
 u=\frac{\sqrt\alpha-b}{\pi}
\]
satisfies the Eisenstein polynomial
\[
 u^2+\frac{2b}{\pi}u+\frac2\pi=0.
\]
It is a uniformizer of $k(\sqrt\alpha)$ and gives an integral power
basis. In its derivative, $2u$ has valuation five and $2b/\pi$ has
valuation two. The relative different has exponent two, and the quadratic
conductor--discriminant formula gives the same exponent for the
character over $k$. Induction to
$\mathbb Q_2$ adds $v_2(\operatorname{disc}k)=2$, so the associated
degree-two representation $\rho$ has conductor exponent four.
For the twist by $\chi_2$, dividing $2\alpha$ by the square $\pi^2$
gives $\alpha'=2-i$. The element $u'=\sqrt{\alpha'}-1$ has Eisenstein
polynomial
\[
 (u')^2+2u'+1-\alpha'=0,
\]
since $v_k(1-\alpha')=1$. Its derivative $2(u'+1)$ has valuation four
in the quadratic extension. Thus the relative conductor is four and
$\rho\otimes\chi_2$ has conductor exponent six. These calculations
use the integral-basis formula for the different and the
conductor--discriminant formula for a quadratic character.

The local genus field $E_2=\mathbb Q_2(i,\sqrt2,\sqrt5)$ has
discriminant exponent sixteen. The conductor--discriminant identity for
$A(\sqrt2)$ gives
\[
 v_2(\operatorname{disc}(A(\sqrt2)))=16+2(4+6)=36.
\]
Its normalized different is $36/16=9/4$, which is unchanged by the
unramified base extension to $L_2$.
\end{proof}

For the choice of generators in this lemma, the elementary reciprocity
classes are $5,-1,-2$, in that order. Here reciprocity classes refer to the
maximal elementary abelian quotient. The characters of $-1$ and $5$
on these generators are respectively $y+z$ and $z$. Consequently the
central square form of $A$ is $(y+z)z$, while the unramified quartic
contributes $z^2$. Their independent sum gives the two central directions
$[y,z]$ and $z^2$ in the presentation. Inertia is
$\langle x,y,[y,z]\rangle$. In particular, we may choose the indicated
representatives of the two inertia directions to have order two.

The maximal pro-$2$ group of $\mathbb Q_2$ has three generators and
one relation. Local Kummer theory gives generator rank three.
The pro-$2$ cohomology comparison in the global proof injects its
second cohomology into $\operatorname{Br}(\mathbb Q_2)[2]$, and the
nonzero local Hilbert pairing makes this injection an isomorphism.
The Kummer basis dual to $5,-1,-2$ under that pairing is $2,-5,5$.
Its Hilbert matrix is
\[
 \begin{pmatrix}0&1&1\\1&1&0\\1&0&0\end{pmatrix},
\]
as follows, for odd $u,v$, from
\[
 (2^a u,2^b v)_2
 =(-1)^{(u-1)(v-1)/4+a(v^2-1)/8+b(u^2-1)/8}.
\]
By the relation--cup identity \cite[Proposition~3.2]{Quadrelli2021},
an actual defining relation in these generators has quadratic initial
$y^{[2]}+[x,y]+[x,z]$.

\subsection{The Global Cuts and Their Finite Quotient}
Define $G$ from $\mathcal G$ by imposing the full kernel of the local
map at two to $D$, inertia squares at $3,5,7,11,13$, Frobenius squares
at $3,5$, and Frobenius fourth powers at
$7,11,13,17,19,23,29,31$. At an odd ramified prime choose the Frobenius
lift with zero coordinate on its own genus character. We may make this
choice by multiplying by inertia. After the inertia square is imposed, the
local tame relation makes inertia and Frobenius commute, so the specified
even-power cut is independent of this remaining choice. Every cut
vanishes in $E$, so the resulting extension still contains $E$.

For the finite binary calculations, number the genus basis
$(-1,2,3,5,7,11,13)$ by $0,\ldots,6$ and encode
$v=\sum v_ie_i$ by the integer $\sum 2^iv_i$. For odd $p$, let
\[
 f_p=\sum_{i:a_i\ne p}\frac{1-(a_i/p)}2e_i,
 \qquad (a_0,\ldots,a_6)=(-1,2,3,5,7,11,13),
\]
where $(\cdot/p)$ is the Legendre symbol. The relevant vectors are
\[
\begin{array}{c|rrrrrrrrrr}
 p&3&5&7&11&13&17&19&23&29&31\\\hline
 f_p&43&86&77&83&58&60&71&57&38&101.
\end{array}
\]
The inertia vector at $p\in T$ is $e_{i(p)}$, and the dyadic vectors are
$x=2$, $y=53$, $z=89$. Let $S(v)$ and $B(v,w)$ denote the free
restricted square and bracket in degree two. In the ordered basis
\[
 e_0^{[2]},\ldots,e_6^{[2]},\quad [e_i,e_j]\ (0\le i<j\le6)
\]
with lexicographically ordered pairs,
\[
 S(v)=\sum_i v_i e_i^{[2]}+\sum_{i<j}v_iv_j[e_i,e_j],\qquad
 B(v,w)=S(v+w)+S(v)+S(w).
\]
The six original quadratic initials are
\[
 S(e_0),\qquad
 B(e_{i(p)},f_p)+\frac{p-1}{2}S(e_{i(p)})\quad(p\in T),
\]
with scalars reduced modulo two. These formulas follow by expanding the tame
relators and the infinity
square. We use them to compute the first three layers of the quotient.

\begin{lemma}\label{tw:new-retained-quotient}
The actual quotient $Q=G/D_4G$ has order $2^{38}$. It retains the
entire local group $D$, local groups $C_2^2$ at $3,5$,
$C_2\times C_4$ at $7,11,13$, and $C_4$ at
$17,19,23,29,31$. The conjugacy class of actual complex conjugation
in $Q$ has size $2^{12}$. Its elementary image is outside every one
of these decomposition groups' elementary spans.
\end{lemma}
\begin{proof}
Use the genus basis $(-1,2,3,5,7,11,13)$. The dyadic generators have
elementary vectors
\[
 x=2,\qquad y=53,\qquad z=89.
\]
The six original odd and infinite quadratic relations are independent,
and their sum is $y^{[2]}+[x,y]+[x,z]$.
Lemma~\ref{tw:complete-global-presentation} shows that the corresponding
original relators give a complete presentation.

The local presentation in Lemma~\ref{tw:new-local-two} has four
independent quadratic initials. In its local kernel relation quotient,
the actual defining relation $r$ of the maximal pro-$2$ group of
$\mathbb Q_2$ has coefficient one on $y^2$: its quadratic initial is
$y^{[2]}+[x,y]+[x,z]$. The other six local kernel generators have no
$y^{[2]}$ coordinate. Replacing $y^2$ by $r$ preserves generation of
the relation quotient, hence closed normal generation by the pro-$2$
Nakayama argument. Globally $r$ already belongs to the closed normal
closure of the six original relators. We need not impose it
again.

The resulting complete presentation has sixteen independent quadratic
initials: the six originals, five odd inertia squares,
$x^{[2]},[x,y],[x,z]$, and the two Frobenius squares at $3,5$.
It also has the actual cubic $[[y,z],z]$ and ten degree-four relators,
including both $z^4$ and $[y,z]^2$. Reduction of the local square and
commutator forms against these sixteen rows retains exactly the two
independent local classes $z^{[2]}$ and $[y,z]$. The three local
degree-one generators are independent. Thus $G/D_3G$, of order
$2^{7+12}$, already retains all of $D$. Reduction against the same rows
verifies the other displayed local
groups and the separation of the elementary image of complex conjugation.

The free restricted cubic layer has dimension $112$. Bracketing the
sixteen quadratic initials with the seven generators gives rank $92$.
The actual cubic $[[y,z],z]$ is independent and raises this rank to
$93$. To see that there are no other degree-three initials, work in the free quotient
$\mathcal F/D_4\mathcal F$. Its second Zassenhaus subgroup is elementary
abelian and its third is central. The normal closure of the actual
quadratic words $\rho_i$ and actual cubic $\kappa$ is the vector space
spanned by $\rho_i,[\rho_i,X_k],\kappa$. Independence of the sixteen
quadratic initials implies that a combination of the $\rho_i$ can lie in
degree three only with every coefficient zero. Its intersection with
degree three is precisely the span of the displayed brackets
and $\kappa$. Their initials depend only on the quadratic initials, so
unknown cubic coefficients of the actual quadratic words do not change
the calculation. Consequently
\[
 \dim D_3G/D_4G=19,\qquad |Q|=2^{7+12+19}=2^{38}.
\]
In the preceding 28-coordinate convention, the sixteen rows have
hexadecimal masks
\[
\begin{gathered}
\texttt{1},\ \texttt{142104},\ \texttt{1444000},\ \texttt{4480410},\
\texttt{a010820},\ \texttt{d020000},\\
\texttt{4},\ \texttt{8},\ \texttt{10},\ \texttt{20},\ \texttt{40},\
\texttt{2},\ \texttt{1a080},\ \texttt{2c080},\
\texttt{814aab},\ \texttt{42aa056}.
\end{gathered}
\]
We compute the ranks by binary elimination in a free associative algebra.
Embed the free restricted degree-two space in the free associative algebra
on $X_0,\ldots,X_6$ by
$e_i^{[2]}\mapsto X_i^2$ and
$[e_i,e_j]\mapsto X_iX_j+X_jX_i$.
Represent homogeneous cubics by their 343 coefficients on the words
$X_iX_jX_k$. For each of the sixteen rows $r$ and each $k$, insert
$rX_k+X_kr$ into a binary row-echelon basis. This produces rank 92.
Insert the tensor for $[[y,z],z]$ to obtain rank 93.
Running the same operation on all 28 free quadratic basis elements gives
rank 112. Ordinary binary elimination in degree two gives rank 16 and
verifies independence of the two retained local classes. Reduction of
$B(e_0,e_k)$ modulo the sixteen rows gives rank six and kernel
$\langle e_0\rangle$. Reduction of $rX_0+X_0r$, for the 28 free quadratic
basis elements $r$, modulo the cubic relation space gives rank six.
Each calculation uses only the displayed masks and these tensor
operations. The supplementary code performs the same integer operations.

Commutation with $c$ has rank six in the second layer and rank six on
the second layer with values in the third. Its degree-one kernel is
exactly $\langle c\rangle$. Since $D_2Q$ has order $2^{31}$ and its
commutator image with $c$ has order $2^6$, the intersection of the
centralizer with $D_2Q$ has order $2^{25}$. Its image in $Q/D_2Q$ is
the line through $c$, giving $|C_Q(c)|=2^{26}$ and the asserted
conjugacy-class size.
\end{proof}

\subsection{Filtered Fox Images}
We use left Fox coefficients \cite{Fox1953}:
\[
 w-1=\sum_i\partial_iw\,(x_i-1),\qquad
 \partial_i(uv)=\partial_i u+u\partial_i v.
\]
Derivative basis vectors have degree one. If a finite-dimensional space
$M$ has a decreasing filtration $F^0M=M$, eventually zero, write
$H_M(t)=\sum_n\dim(F^nM/F^{n+1}M)t^n$. Subspaces have the induced
filtration and quotients the quotient filtration. For $0<t<1$,
\[
 H_M(t)=(1-t)\sum_{n\ge1}\dim(M/F^nM)t^{n-1}.
\]
Consequently a filtration-preserving surjection $M\to N$ satisfies
$H_N(t)\le H_M(t)$, even without strictness. In particular, for subspaces
$U,V$ of a filtered ambient space,
\begin{equation}\label{tw:filtered-sum}
 H_{U+V}(t)\le H_U(t)+H_V(t)-H_{U\cap V}(t).
\end{equation}
To obtain this inequality, apply additivity to the quotient filtration
from $U\oplus V$, then compare it with the induced filtration on the sum.
We next compute the cost of a full local relation kernel.

\begin{lemma}\label{tw:filtered-local-freeness}
Let $P$ be a finite $2$-group and $D'\le P$ one of
$C_2,C_4,C_2^2,C_2\times C_4$, or the group $D$ of
Lemma~\ref{tw:new-local-two}. Suppose its degree-one directions, and
its nonzero degree-two directions when present, embed in the corresponding
Zassenhaus layers of $P$. If $D'$ has $r$ minimal generators and
augmentation Hilbert polynomial $P_{D'}$, its full induced local Fox
kernel has Hilbert polynomial
\begin{equation}\label{tw:local-full-cost}
 \left(rt-1+P_{D'}(t)^{-1}\right)H_{\mathbb F_2[P]}(t).
\end{equation}
Its image under a filtration-preserving map has no larger Hilbert value.
\end{lemma}
\begin{proof}
These local groups have no nonzero layer above degree two, so the
hypotheses imply $D'\cap D_nP=D_nD'$ for every $n$.
Extend a homogeneous basis of the graded restricted Lie algebra of $D'$
to one for $P$, ordering the local factors last. The restricted PBW
form of Jennings' theorem \cite{Jennings,Quillen} gives lifts $t_j$ of the
remaining monomials such that
\[
 \mathbb F_2[P]=\bigoplus_jt_j\mathbb F_2[D']
\]
as right modules and filtered vector spaces. The shift of the $j^{th}$
summand is the weight of $t_j$, and the shift polynomial is
$H_{\mathbb F_2[P]}/P_{D'}$.
Put $B=\mathbb F_2[D']$. The local map
$B^r(-1)\to I_{D'}$, $(a_i)\mapsto\sum a_i(d_i-1)$, is strict, because
$I_{D'}^n=\sum_iI_{D'}^{n-1}(d_i-1)$. Its kernel has polynomial
$rtP_{D'}-(P_{D'}-1)$. Filtered right-freeness preserves exactness and
strictness under induction, proving the formula. Finally, the local
generator words have independent elementary images in a minimal global
free presentation and so extend to a free pro-$2$ basis. The Fox
change-of-basis matrix and its inverse preserve augmentation degree.
Thus the identification with a coordinate submodule in the global
derivative module preserves the induced filtration. The assertion about
images follows from the cumulative Hilbert formula.
\end{proof}

The dyadic block shares a relation with the odd and infinite blocks.
We compute the image of this relation as well as the full dyadic cost.
\begin{lemma}\label{tw:new-local-fox}
For this dyadic local group the full induced local Fox-image cost is
$d_D(t)=3t-1+P_D(t)^{-1}$. The common actual dyadic defining row has
normalized image polynomial
\[
 s_D(t)=t^2\left(1-\frac{t^7}{P_D(t)}\right).
\]
\end{lemma}
\begin{proof}
The independent local degree-one and degree-two classes in
Lemma~\ref{tw:new-retained-quotient} give an embedding of the local
restricted graded Lie algebra in the global one. Extend its restricted
PBW basis to a global basis, with the local factors last. This proves
filtered freeness as a right local-algebra module, including the exact
augmentation shifts. The local augmentation sequence then gives the
full Fox-image cost $3t-1+P_D^{-1}$ by
Lemma~\ref{tw:filtered-local-freeness}. This argument does not require
commutativity.

Write $B=\mathbb F_2[D]$. The linear coefficients of the actual
local relation's Fox row span $I_B/I_B^2$, because the local Hilbert
pairing is nonsingular \cite[Chapter~III, Theorem~4.4]{MilneCFT}, and
its matrix is the matrix of linear Fox coefficients
\cite[Proposition~3.2]{Quadrelli2021}. In our coordinates the relation
has initial $y^{[2]}+[x,y]+[x,z]$, whose three derivative coefficients
are $y+z,y+x,x$ and are independent. They generate $I_B$ as a right ideal by
Nakayama. Their common left annihilator is the one-dimensional left socle of $B$.
We can check the required graded assertion in this 32-element algebra.
Starting with its displayed normal-form multiplication,
form $J^0=B$ and obtain $J^{n+1}$ by spanning the right products of a
basis of $J^n$ by $x-1,y-1,z-1$. The dimensions for $0\le n\le8$ are
\[
 32,31,28,23,16,9,4,1,0.
\]
On $J^n/J^{n+1}$ the simultaneous multiplication map to
$(J^{n+1}/J^{n+2})^3$ has kernel zero for $n<7$ and dimension one for
$n=7$. This is a finite row reduction on vectors of length 32 using the
specified multiplication. The top vector is the sum of all group elements.
Changing the three degree-one coefficients by an invertible matrix does
not change these kernels. Thus every nonzero homogeneous element below
degree seven has nonzero product with at least one linear Fox coefficient. Row
multiplication raises coefficient degree by exactly one, and the
derivative basis shift contributes one more. The row image consequently has
polynomial $t^2(P_D-t^7)$. Filtered induction gives the stated normalized
cost.

The actual row belongs to the full dyadic block because the local
relation vanishes in $D$. It also belongs to the sum of the odd blocks
and infinity square. Indeed, by the complete global presentation, the
local relation belongs to the closed normal closure of those original
relators. If a defining relator $\rho$ evaluates to one in the finite
group, the Fox product rule gives
\[
 \partial(u\rho u^{-1})=u\partial\rho.
\]
Products of conjugates give module combinations of the original rows.
Evaluation in the finite derivative module is continuous
and its row span is closed, so the same conclusion holds for the closed
normal closure. The filtered local coordinate change used above preserves
the computed image polynomial. Thus the common submodule is an actual
intersection in the global derivative module.
\end{proof}

\subsection{Infinitude and the Resulting Number Fields}
The local Fox bounds imply that the global group is infinite and give
the following family of number fields.
\begin{theorem}\label{tw:new-field-family}\label{tw:field-family}
The group $G$ is infinite. For every sufficiently large power of two $d$
there is a totally imaginary Galois extension $K/\mathbb Q$ of degree
$2d$ retaining $Q=G/D_4G$. For its actual complex conjugation $c$, the
fixed field $F=K^{\langle c\rangle}$ has at least one real embedding and
\[
 \operatorname{rd}(K)=\operatorname{rd}(F)
 =2^{9/4}\sqrt{15015},\qquad
 \frac{4095}{8192}\le\theta:=\frac{r_2(F)}d<\frac12.
\]
The extension $K/F$ is quadratic, unramified at every finite place, and
split at every place above the following rational primes, whose uniform
local types in both $K$ and $F$ are
\[
\begin{array}{c|c|c}
 p&e_p&f_p\\\hline
 2&8&4\\
 3,5&2&2\\
 7,11,13&2&4\\
 17,19,23,29,31&1&4.
\end{array}
\]
\end{theorem}
\begin{proof}
Suppose first that $G$ is finite and put $A=\mathbb F_2[G]$.
The strict global augmentation map $A^7(-1)\to I_G$ has kernel $L$ with
$H_L=1-(1-7t)H_A$. The continuous Fox relation sequence says that the
rows of the complete defining relations generate $L$. This also follows
by tensoring the free augmentation sequence and using the ideal normally
generated by the defining relators.
The local hypotheses of Lemma~\ref{tw:filtered-local-freeness} hold by
Lemma~\ref{tw:new-retained-quotient}. Write
\[
 c_1=\frac{t^2}{1+t},\qquad c_2=2t-1+(1+t)^{-2},\qquad
 c_{24}=2t-1+\frac1{(1+t)^2(1+t^2)},\qquad
 c_4=\frac{t^4}{(1+t)(1+t^2)}.
\]
The two odd $C_2^2$ blocks cost $2c_2H_A$, the three odd
$C_2\times C_4$ blocks cost $3c_{24}H_A$, infinity costs $c_1H_A$,
and the five unramified order-four caps cost $5c_4H_A$.
The full dyadic block costs $d_DH_A$. Its actual intersection with the
sum of the odd and infinite blocks contains the common row image
$s_DH_A$, which is subtracted once by \eqref{tw:filtered-sum}.
Consequently
\[
 1\le H_A(t)\mathcal P(t),\qquad
 \mathcal P(t)=1-7t+2c_2+3c_{24}+d_D+c_1-s_D+5c_4.
\]
Exact rational substitution gives
\[
 \mathcal P(11/34)=-\frac{160218343}{112275691650}<0.
\]
This contradicts finiteness. The inequality includes every defining
relation, and the subtracted term comes from the common row image in the
global Fox kernel.

An infinite profinite group has finite quotients of unbounded order.
Taking common refinements with $Q$ gives arbitrarily large finite
$2$-quotients $H\twoheadrightarrow Q$. A nontrivial kernel meets $Z(H)$
nontrivially by the class equation for conjugation, and hence contains a
central subgroup of order two. Repeated quotienting by such subgroups
inside the kernel realizes every intermediate power-of-two order while
retaining $Q$. These quotients correspond to the claimed Galois fields.

Because $K$ contains $i$, it is totally imaginary. The chosen complex
conjugation fixes a real subfield in the given embedding, so $F$ has a
real embedding. If $b=r_1(F)$ and $H=\operatorname{Gal}(K/\mathbb Q)$,
embeddings of $F$ correspond to cosets of $\langle c\rangle$ in $H$.
A coset is real exactly when its representative centralizes $c$. Hence
\[
 b=|C_H(c)|/2,\qquad \frac bd=\frac{|C_H(c)|}{|H|}.
\]
The conjugacy class in $H$ maps onto the class in $Q$, of size $2^{12}$.
Thus $0<b/d\le2^{-12}$, giving the stated interval for
$\theta=(1-b/d)/2$.

Conjugation preserves an elementary image. The separation verified in
Lemma~\ref{tw:new-retained-quotient} excludes every conjugate
of $c$ from each ramified or selected decomposition group. Intersection
with $\langle c\rangle$ is trivial, proving relative splitting at these
places. Elsewhere $K/\mathbb Q$ is unramified, so $K/F$ is unramified
there also. Relative splitting transfers the displayed local types to $F$.
Tame inertia of order two contributes $p^{1/2}$ to the root discriminant
at each $p\in T$, and the dyadic lemma contributes $2^{9/4}$. Finally
$\Delta_K=\Delta_F^2$ because the relative finite discriminant is one.
\end{proof}

\section{An Explicit Retained Field and Its Hecke Factors}\label{ar:field}
The analytic estimate uses a fixed subfield of every sufficiently large
field in Theorem~\ref{tw:field-family}. We construct this subfield by
Kummer theory and express its zeta function as a product of Dirichlet and
Hecke functions. Exact local calculations then determine the conductors
and Euler polynomials needed to evaluate these factors.

\subsection{A Catalog of Norm Extensions}\label{ar:catalog}
Retain $V=\operatorname{Gal}(E/\mathbb Q)=\mathbb F_2^7$ with the
coordinates of Section~\ref{tw:tower}. For a rational squareclass $A$
supported on $-1,2,3,5,7,11,13$, write $\chi_A\in V^*$ for its Kummer
character. Define $\alpha_i=x_i+y_i\sqrt{A_i}$ by the following seventeen rows.
Each row satisfies
$N_{\mathbb Q(\sqrt{A_i})/\mathbb Q}(\alpha_i)=B_i z_i^2$.
\begin{equation}\label{ar:norm-table}
\begin{array}{r|r|r|r|r|r}
 i&A_i&B_i&x_i&y_i&z_i\\\hline
0&-1&65&4&7&1\\
1&-55&91&6&1&1\\
2&-2&33&1&4&1\\
3&-55&14&1&1&2\\
4&-55&26&7&1&2\\
5&-7&2&1&1&2\\
6&-10&1001&1&10&1\\
7&22&273&19&2&1\\
8&-39&55&4&1&1\\
9&-35&429&1&7&2\\
10&-143&42&5&1&2\\
11&-3&13&1&2&1\\
12&-26&105&1&2&1\\
13&-6&385&1&8&1\\
14&-14&65&3&2&1\\
15&-11&5&3&1&2\\
16&-1&13&-2&3&1
\end{array}
\end{equation}
All these numbers are integral. Their base discriminants and norms are
supported on $\Sigma=\{2,3,5,7,11,13\}$, so their Galois closures are
unramified outside $\Sigma$. An individual catalog field need not survive
all the cuts in $G$. We will check retention for the products used in
the construction.

For a nonnegative integer $a<2^{17}$ define
\[
 \beta(a)=\prod_{i=0}^{16}\alpha_i^{a_i},\qquad
 q_a(v)=\sum_{i=0}^{16}a_i\chi_{A_i}(v)\chi_{B_i}(v),
 \quad a=\sum_i2^ia_i.
\]
Products and signs in the second formula are taken in $\mathbb F_2$.
The norm identities determine the central extension associated to each
product.
\begin{lemma}\label{ar:kummer}
For any catalog word $a$, the field $E(\sqrt{\beta(a)})$ is Galois over
$\mathbb Q$, and its invariant Kummer class has transgression $q_a$. When this form
is nonzero, the field is a central extension of $V$ by $C_2$ with
that quadratic form. Products of the catalog radicands add their extension classes.
On the space of invariant Kummer classes in the maximal pro-$2$
$\Sigma$-unramified extension above $E$, this class map is injective.
\end{lemma}
\begin{proof}
Choose the square roots of the positive rational numbers $B_i$ compatibly
inside $E$. If $v\in V$ changes the sign of $\sqrt{A_i}$, put
\[
 u_{i,v}=\frac{z_i\sqrt{B_i}}{\alpha_i}.
\]
Otherwise put $u_{i,v}=1$. The norm identity gives
\[
 \alpha_i u_{i,v}^2=v(\alpha_i),\qquad
 u_{i,v}v(u_{i,v})=(-1)^{\chi_{A_i}(v)\chi_{B_i}(v)}.
\]
Thus $v$ lifts to an automorphism by sending
$\sqrt{\alpha_i}$ to $u_{i,v}\sqrt{\alpha_i}$. For a product radicand
use $u_v=\prod_i u_{i,v}^{a_i}$. It satisfies
\begin{equation}\label{ar:kummer-lifts}
 \beta(a)u_v^2=v(\beta(a)),\qquad
 u_vv(u_v)=(-1)^{q_a(v)}.
\end{equation}
The kernel over $E$ is central of exponent two. The squares of lifted
elements are given by $q_a$, and their commutators by its
polar form $B_{q_a}(v,w)=q_a(v+w)+q_a(v)+q_a(w)$. These squares and
commutators determine the central extension, since every word can be
collected into an ordered product of the seven lifted basis elements and
a central sign. Applying the same calculation to a product of radicands
gives the addition rule.

For injectivity let $\mathcal H$ be the subgroup of $\mathcal G$ fixing
$E$. The five-term sequence contains
\[
 H^1(\mathcal G,\mathbb F_2)\longrightarrow
 H^1(\mathcal H,\mathbb F_2)^V
 \xrightarrow{\mathrm{trg}} H^2(V,\mathbb F_2).
\]
Every quadratic character of $\mathcal G$ factors through $E$, by the
Kummer computation in Lemma~\ref{tw:complete-global-presentation}.
The first arrow is zero, so transgression is injective.
The calculation of squares and commutators above identifies this
transgression with the indicated quadratic form. In particular a
nonzero form cannot come from a square radicand in $E$, and independent
forms give independent Kummer classes.
\end{proof}

\subsection{The Completed Field and the Full Quadratic Quotient}
Define two lists of raw catalog words
\begin{align}
 \mathcal A&=(17,27184,6220,2,32768,39508,132),
       \label{ar:completed-words}\\
 \mathcal B&=(2,256,512,4096,9,65,132,2052,24576,37,32768,65536).
       \label{ar:full-words}
\end{align}
Set
\[
 N=E\bigl(\sqrt{\beta(a)}:a\in\mathcal A\bigr),\qquad
 M=E\bigl(\sqrt{\beta(b)}:b\in\mathcal B\bigr).
\]
For example, writing $\eta=\alpha_{13}\alpha_{14}$, the seven defining
radicands of $N$ are
\begin{align*}
 \gamma_1&=\alpha_0\alpha_4,&
 \gamma_2&=\alpha_4\alpha_5\alpha_9\alpha_{11}\eta,\\
 \gamma_3&=\alpha_2\alpha_3\alpha_6\alpha_{11}\alpha_{12},&
 \gamma_4&=\alpha_1,\\
 \gamma_5&=3+\sqrt{-11},&
 \gamma_6&=\alpha_2\alpha_4\alpha_6\alpha_9\alpha_{11}\alpha_{12}\gamma_5,\\
 \gamma_7&=(1+4\sqrt{-2})(19+2\sqrt{22}).
\end{align*}

These products give two subfields of the tower.
\begin{proposition}\label{ar:central-fields}
The fields $N$ and $M$ are contained in every field $K$ of
Theorem~\ref{tw:field-family}. Their degrees are
\[
 [N:\mathbb Q]=2^{14}=16384,\qquad [M:\mathbb Q]=2^{19},\qquad N\subset M.
\]
The field $M$ is exactly the field corresponding to $G/D_3G$.
\end{proposition}
\begin{proof}
Use the 28-coordinate convention dual to the restricted quadratic basis
in Section~\ref{tw:tower}: the first seven coordinates encode $v_i^2$
and the remaining 21 encode $v_iv_j$ for lexicographically ordered
$i<j$. In this convention the forms for $\mathcal A$ and $\mathcal B$
have hexadecimal masks
\begin{equation}\label{ar:central-masks}
\begin{array}{c|l}
\mathcal A&
\texttt{9014280},\ \texttt{af24500},\ \texttt{6830680},\ \texttt{b401400},\\
&\texttt{800200},\ \texttt{acf0980},\ \texttt{a138900}\\[2pt]
\mathcal B&
\texttt{b401400},\ \texttt{9140a00},\ \texttt{78c1900},\ \texttt{520e700},\\
&\texttt{2415680},\ \texttt{1c38e00},\ \texttt{a138900},\ \texttt{213900},\\
&\texttt{45f9c00},\ \texttt{1bb80},\ \texttt{800200},\ \texttt{1000}.
\end{array}
\end{equation}
To obtain these masks, expand each product of linear characters from
\eqref{ar:norm-table}, writing a diagonal term as $v_i^2$.
Binary elimination gives ranks seven and twelve, respectively, and
shows that the first span is contained in the second. Taking the binary
scalar product with each of the sixteen displayed relation rows in
Section~\ref{tw:tower} gives zero for every form in the second list.
The rows have rank 16 in dimension 28, so these twelve independent forms
span their entire annihilator.

Lemma~\ref{ar:kummer} now proves the asserted Kummer degrees and the
containment. Each extension has central kernel of exponent two over the
exponent-two group $V$, and hence class at most two and exponent at most
four. Its third Zassenhaus subgroup is thus trivial. Evaluating a defining
relator of $G$ in this group depends only on its quadratic initial, and
the higher-degree relators vanish. Since the forms annihilate every
quadratic relation, these number fields satisfy all the defining
relations of $G$, including the local cuts. Finally
$G/D_3G$ has the same order $2^{19}$ as $M$, by
Lemma~\ref{tw:new-retained-quotient}, which proves equality. Since the
fields $K$ retain $G/D_4G$, they retain both $M$ and $N$.
\end{proof}

\subsection{The Complete Representation Inventory}
For a quadratic form $q$ on $V$, write $R_q=\operatorname{rad}(B_q)$
and $g(q)=\sum_{v\in V}(-1)^{q(v)}$. The space of forms defining $N$
will be denoted $\mathcal Q=\langle q_a:a\in\mathcal A\rangle$.
To count the irreducible representations and their gamma factors, we
examine the polar forms and the action of complex conjugation.
\begin{lemma}\label{ar:representations}
The irreducible representations of $\operatorname{Gal}(N/\mathbb Q)$
comprise 128 of degree one, 256 of degree two, 456 of degree four, and
124 of degree eight. Exactly eight of the degree-four representations
are pure at infinity: four have gamma vector $(0,0,0,0)$ and four have
$(1,1,1,1)$. Every other nonlinear representation is mixed, with equal
numbers of zeros and ones in its gamma vector.
\end{lemma}
\begin{proof}
Fix a nontrivial character of $\operatorname{Gal}(N/E)$ and let
$q\in\mathcal Q\setminus\{0\}$ be its central form. Suppose
$\operatorname{rank}B_q=2m$. The corresponding central pushout has
commutator pairing $B_q$, and its center above $V$ is the inverse image
of $R_q$. There are $|R_q|=128/2^{2m}$ central characters extending the
chosen nontrivial kernel character. Each gives one irreducible
representation of dimension $2^m$.
To construct these representations, choose a maximal isotropic subspace
of $V$ for $B_q$, of dimension $7-m$ and containing $R_q$. Its inverse image is abelian.
Extend the chosen central character to it and induce to the full group.
Nondegeneracy of the induced pairing on $V/R_q$ shows, by the usual
intertwiner criterion for induction, that no element outside that
subgroup intertwines its character. The induced representation is
irreducible and has degree $2^m$. Changing the extension inside this
subgroup only conjugates the character, while changing its restriction
to the center gives distinct representations. Their squared dimensions
sum to 128, so this accounts for the entire central sector.
The trivial central character gives the 128 linear characters of $E$.

Enumeration of the 127 nonzero sums of the seven masks in
\eqref{ar:central-masks} gives
\[
\begin{array}{c|rrr}
 \operatorname{rank}B_q&2&4&6\\\hline
 \#q&8&57&62\\
 \#\{q:g(q)>0\}&5&34&48\\
 \#\{q:g(q)=0\}&3&23&14.
\end{array}
\]
The positive sums are respectively 64, 32, and 16. No negative Gauss
sum occurs. For this enumeration, evaluate each quadratic polynomial on the 128
binary vectors and compute the rank of its $7$ by $7$ polar matrix.
The resulting dimension count is
\[
 128+256\cdot2^2+456\cdot4^2+124\cdot8^2=16384.
\]

The linear map $q\mapsto B_q(c,\cdot)$ on $\mathcal Q$, where
$c=e_0$, has rank six. Its only nonzero kernel form has rank four
and is represented by $\alpha_4\alpha_6\alpha_9\alpha_{10}\alpha_{11}$.
Every form satisfies $q(c)=0$, as it annihilates the infinity square.
If $B_q(c,\cdot)\ne0$, an element conjugating $c$ to its product with
the nontrivial central sign shows that the trace of $c$ in the
representation is zero. Because $c^2=1$, its $+1$ and $-1$ eigenspaces
then have equal dimensions. For the unique remaining sector $c$ is
central in the pushout, hence acts scalarly. Its eight genus twists
split equally between the two signs. The explicit description below
identifies the even and odd twists. The dimensions of these eigenspaces
give the asserted real gamma factors.
\end{proof}

The regular representation gives the exact Artin factorization
\begin{equation}\label{ar:artin-factorization}
 \zeta_N(s)=\prod_{\rho\in\operatorname{Irr}(\operatorname{Gal}(N/\mathbb Q))}
                     L(s,\rho)^{\dim\rho}.
\end{equation}
This follows prime by prime by decomposing the regular representation
and taking the determinant on inertia invariants. We next identify every
nonlinear factor with an entire Hecke function,
which allows us to use this identity without an Artin holomorphy
assumption.

\subsection{Descent to Quadratic Hecke Characters}\label{ar:descent}
We construct the inducing characters by descending the catalog radicands
to suitable subfields of $E$. The descent gives exact identities for the
Artin functions in the analytic calculation.

If $g(q)>0$ and $\operatorname{rank}B_q=2m$, then $q$ vanishes on its
radical and the induced nondegenerate quadratic form has Arf invariant
zero. Equivalently it has a totally singular subspace $U$ of dimension
$7-m$ containing $R_q$. We may find such a subspace by enumerating the subspaces of $V$ and
testing $q(u)=0$ on each vector. Put $B=E^U$, so
$[B:\mathbb Q]=2^m$, and choose a catalog product $\beta$ representing $q$.
Let $u_v$ be the exact lifts from \eqref{ar:kummer-lifts}. On $E$ define
$T_vP=u_vv(P)$. For $v\in U$ these are commuting involutions: their
squares and commutators are respectively $(-1)^{q(v)}$ and
$(-1)^{B_q(v,w)}$.

Starting from $P=1$, take a basis $v_1,\ldots,v_{7-m}$ of $U$ and
successively replace
\begin{equation}\label{ar:projectors}
 P\longmapsto P+T_{v_j}P
 \quad\hbox{or, if this is zero,}\quad P-T_{v_j}P.
\end{equation}
At least one choice is nonzero in characteristic zero. Commutativity
preserves the eigenspaces already selected, so the final nonzero $P$
satisfies $T_vP=\pm P$ for every $v\in U$. Consequently
\[
 \eta=\beta P^2\in B,
 \qquad E(\sqrt\eta)=E(\sqrt\beta).
\]
Every operation is rational arithmetic in the basis
$\prod_i\sqrt{a_i}^{\epsilon_i}$ of $E$. At a step, if
$\eta_0=\beta P^2$, the element
$h=\beta u_vP v(P)$ satisfies
\[
 h^2=\eta_0v(\eta_0),\qquad
 \beta(P\pm T_vP)^2=\eta_0+v(\eta_0)\pm2h.
\]
At each descent step, these identities verify the new radicand, while
the nonzero test and the invariance under the processed subgroup verify
the chosen projector.

Let $L=B(\sqrt\eta)$ and let $\lambda$ be its nontrivial quadratic
Hecke character over $B$. The inverse image of $U$ in the central
pushout is the elementary abelian subgroup used in the representation
proof. Its character selecting $\sqrt\eta$ induces the required
irreducible representation. Thus
\begin{equation}\label{ar:quadratic-induction}
 L(s,\rho)=L_B(s,\lambda)=\frac{\zeta_L(s)}{\zeta_B(s)}.
\end{equation}
This is an entire nontrivial Hecke function. Its completed functional
equation has root number $+1$, because the completed Dedekind zeta
functions in the quotient have root number $+1$.
The remaining representations in this sector are obtained by multiplying
$\eta$ by rational squareclasses $D$ representing
$V^*/B_q(V)$. There are $|R_q|$ such twists. Here $B_q(V)$ denotes the
image of $v\mapsto B_q(v,\cdot)$. Twisting by a character in this image
conjugates the inducing representation, while the quotient distinguishes
its central characters.

If $g(q)=0$, choose $p\in\{5,13\}$ so that $g(q+\chi_p)>0$.
For this seven-dimensional space, exact Gauss-sum evaluation gives
26 sectors permitting either choice, nine requiring $p=5$, and five
requiring $p=13$. Thus a suitable choice exists in every sector. Put
\[
 \beta_5=-10-2\sqrt5,\qquad \beta_{13}=-26-6\sqrt{13}.
\]
Their norms are $16p$ and their central forms are $\chi_p$.
To identify these as the cyclic quartic extensions of conductor $p$,
we use their defining polynomials
\[
 f_5(X)=X^4+20X^2+80,\qquad f_{13}(X)=X^4+52X^2+208.
\]
If $r^2=\beta_p$, the conjugation rule $r\mapsto4\sqrt p/r$ has
square $r\mapsto-r$, and gives their cyclic Galois group of order four.
The following integral bases have trace-pairing discriminants $5^3$
and $13^3$, respectively:
\[
\begin{split}
 p=5:\quad &1,\quad r^2/8-r/4+3/2,\quad-r^2/8-r/4-3/2,\quad
 r^3/16-r^2/8+3r/4-1;\\
 p=13:\quad &1,\quad r^2/24-r/4+5/6,\quad-r^2/24-r/4-5/6,\quad
 r^3/48-r^2/24+11r/12-4/3.
\end{split}
\]
Integrality and the discriminants follow by substitution in $f_p$ and
rational trace-matrix arithmetic. These orders are maximal. The field is
unramified outside $p$, and its
quadratic subfield $\mathbb Q(\sqrt p)$ is ramified at $p$. Hence inertia
is the full cyclic group of order four. Tame ramification gives
discriminant exponent three, which agrees with the computed order
discriminant.
The conductor--discriminant formula gives conductor $p$ for both faithful
quartic characters, identifying them with $\xi_p$ and
$\overline{\xi_p}$ under global reciprocity.
Descend $\beta\beta_p$ by the positive construction. If $\rho'$ is
its quadratic induction and $\xi_p$ is the primitive quartic Dirichlet
character normalized by $\xi_p(2)=i$, use $\rho'\otimes\xi_p$.
The two quadratic forms add to
$(q+\chi_p)+\chi_p=q$, so this is a representation of the actual
original central field. This can also be seen directly on the two
Kummer square roots: the extra central sign from $\beta_p$ occurs
in both factors and cancels in their tensor product. Its genus twists
again give every representation in the sector. Since
$\overline{\xi_p}=\xi_p\chi_p$ and $\rho'$ is real, the genus twists
pair the complex representation with its conjugate. The resulting
functions are entire Hecke functions with root numbers of modulus one.
The bounds below do not require their phases to be known.

\subsection{The Pure Quartic Field}
The unique pure form in Lemma~\ref{ar:representations} is represented
by $\beta=\alpha_4\alpha_6\alpha_9\alpha_{10}\alpha_{11}$.
Five projector steps give
\[
 B_0=\mathbb Q(\sqrt{30},\sqrt{65}),\qquad
 \eta_0=393911-10360\sqrt{30}+23013\sqrt{65}-2400\sqrt{1950}.
\]
To verify this squareclass identity, start with $P_0=1$ and the
multipliers $u_v$ for this $\beta$. Apply
$P_j=P_{j-1}+\epsilon_jT_{v_j}P_{j-1}$ with
\[
 (v_1,\ldots,v_5)=(1,16,32,6,74),\qquad
 (\epsilon_1,\ldots,\epsilon_5)=(1,1,1,-1,-1).
\]
Exact multiplication in $E$ gives $\eta_0=\beta P_5^2/256$.
The five vectors span the simultaneous kernels of $\chi_{30}$ and
$\chi_{65}$, so their fixed field is $B_0$. All four real values are
positive: using $\sqrt{30}<6$, $\sqrt{65}<9$, and $\sqrt{1950}<45$,
each is larger than
$393911-10360\cdot6-23013\cdot9-2400\cdot45=16634$. The eight factors are
\begin{equation}\label{ar:pure-factors}
 L_D(s)=\frac{\zeta_{B_0(\sqrt{D\eta_0})}(s)}{\zeta_{B_0}(s)},
 \qquad D\in\{\pm1,\pm2,\pm3,\pm5\}.
\end{equation}
The base discriminant is
$120\cdot65\cdot312=2433600$. For every listed $D$, the maximal-order
discriminant of $B_0(\sqrt{D\eta_0})$ is
\[
 2^{14}3^4 5^4 7^2 11^2 13^4=140455850895360000.
\]
The ratio is $240240^2$, giving the conductor of each factor by the
conductor--discriminant formula. Positive $D$ have four even gamma
factors, negative $D$ four odd ones, and every root number is $+1$.
The discriminants here are those of maximal orders, as required by the
formula below.

\subsection{Exact Local Data}
The inducing characters determine the completed functions used in the
analytic estimate.
\begin{proposition}\label{ar:hecke-data}
All irreducible factors of $\zeta_N$ are the actual Dirichlet factors
and Hecke inductions just constructed. The gamma inventory is the one
in Lemma~\ref{ar:representations}. Their conductors and exceptional
Euler denominators are given by the following finite local algorithms,
with exact inducing-field, radicand, twist, and projector data in the
supplement. In particular no conjectural Artin holomorphy, conductor
identification from a truncated Euler product, or guessed root number
is used.
\end{proposition}
\begin{proof}
The preceding constructions give every irreducible representation in
the regular representation. To compute its Hecke data in a positive
sector with twist $D$, compute the maximal orders of
$B$ and $L_D=B(\sqrt{D\eta})$. The conductor-discriminant formula gives
\begin{equation}\label{ar:conductor-ratio}
 Q_\rho=\frac{|\Delta_{L_D}|}{|\Delta_B|}
       =|\Delta_B|\,N_{B/\mathbb Q}(\mathfrak d_{L_D/B}).
\end{equation}
For a prime $p$, factor $p$ in both maximal orders. With
$f(\mathfrak p/p)$ denoting residue degrees, put
\[
 D_{B,p}(T)=\prod_{\mathfrak p\mid p}(1-T^{f(\mathfrak p/p)}),
 \qquad
 D_{L_D,p}(T)=\prod_{\mathfrak P\mid p}(1-T^{f(\mathfrak P/p)}).
\]
The exact local denominator of \eqref{ar:quadratic-induction} is
\begin{equation}\label{ar:denominator-ratio}
 D_{\rho,p}(T)=D_{L_D,p}(T)/D_{B,p}(T).
\end{equation}
This quotient is a polynomial: above each prime of $B$, quadratic
splitting contributes $1-T^f$, inertia degree two contributes $1+T^f$,
and ramification contributes 1. Its degree is at most $\dim\rho$.
At a real place its character is even or odd according as $D\eta$
is positive or negative. At a complex place the induction contributes
one even and one odd real gamma factor.

We can also compute these data locally, including in the sectors with a
quartic twist. At an odd prime of $B$, first remove the
even valuation of the quadratic radicand by a square. An odd remaining
valuation is ramified with conductor exponent one. For valuation zero,
the residue quadratic character determines splitting, giving the
unramified value of the quadratic character on a uniformizer. At a
dyadic place, use the finite unit quotient and square test: a unit
$1+u$ with valuation $v(u)>2v(2)$ is a square by Hensel's lemma. Thus
squareclass and quadratic-character computations reduce to a finite
quotient of the local integer ring. Its ramification and conductor can
be determined by testing the character on the descending unit groups.
The quadratic norm criterion determines each character value. All these tests
use exact finite-ring arithmetic.

For $\rho=\operatorname{Ind}_{B}^{\mathbb Q}\lambda\otimes\xi_p$, apply
these tests to the local character
\[
 \lambda_{\mathfrak p}\,\bigl(\xi_{p,\ell}\circ
 N_{B_{\mathfrak p}/\mathbb Q_\ell}\bigr)
 \qquad(\mathfrak p\mid\ell).
\]
The local component of a primitive Dirichlet character is normalized to
have the usual Frobenius value away from its modulus. At its modulus,
its restriction to units is the inverse of the finite Dirichlet
character, and its value on the rational uniformizer is one.
For each local character, its conductor exponent is the least integer
$n$ for which it is trivial on $1+\mathfrak p^n$, with exponent zero
for an unramified character. The search is finite: its conductor is no
larger than the maximum of the two known factor conductors. Induction
then gives
\[
 a_\ell(\rho)=\sum_{\mathfrak p\mid\ell}
 f(\mathfrak p/\ell)\bigl(a_{\mathfrak p}(\lambda\xi_p\circ N)
                         +d_{B_{\mathfrak p}/\mathbb Q_\ell}\bigr),
\]
where $d_{B_{\mathfrak p}/\mathbb Q_\ell}$ is the exponent of the local
different. If a local inducing character is ramified it contributes no
Euler factor. If it is unramified, with uniformizer value $\varepsilon$,
it contributes $1-\varepsilon T^{f(\mathfrak p/\ell)}$. Multiplication over
$\mathfrak p\mid\ell$ supplies the target denominator with its actual
complex phase.

At the cancelling odd prime this calculation has a particularly short
form. Write a local uniformizer as $\pi^2=pr$, let $u$ be the normalized
quadratic radicand unit, and write $f$ for the residue degree. The
cancelled character is unramified and its value is
\begin{equation}\label{ar:cancelled-value}
 \varepsilon=(u/\mathbb F_{p^f})\,
       \xi_p((-r)^f)^{-1}.
\end{equation}
The first factor is the value of the unramified quadratic character,
and the second comes from the unit part of the norm of $\pi$. The
inverse is required by the local normalization at the modulus. This
determines each complex Euler polynomial separately.

The finite arithmetic records specify the rational defining data,
projector identities, primitive character values, and exceptional
polynomials. The supplementary checks apply the displayed identities
and local operations to these records. Independent maximal-order and
prime-factorization computations check the discriminant ratios and
local polynomials a second time. The identification of the infinite
Euler products follows from the inducing characters constructed above.
\end{proof}

\subsection{Frobenius Tests for the Prime Budget}\label{ar:frobenius}
For an ordinary prime $p\notin\Sigma$, let $v_p\in V$ be its genus
Frobenius vector. If $v_p\ne0$, a lift in either central field has square
measured by the corresponding quadratic forms. Consequently
\begin{equation}\label{ar:nonsplit-test}
 \begin{aligned}
 f_N(p)=4&\ \Longleftrightarrow\ q_a(v_p)=1\text{ for some }a\in\mathcal A,\\
 f_M(p)=4&\ \Longleftrightarrow\ q_b(v_p)=1\text{ for some }b\in\mathcal B.
 \end{aligned}
\end{equation}
If no form detects the square, the residue degree is two, since the
genus vector is nonzero. In particular a prime with $f_N(p)=2$ and
$f_M(p)=4$ has residue degree at least four in every field $K$.

If $v_p=0$, choose any square root $s_i$ of $A_i$ modulo $p$ and set
\[
 \varepsilon_i(p)=\left(\frac{x_i+y_i s_i}{p}\right)\in\{\pm1\}.
\]
The norm identity shows both that the argument is nonzero and that the
sign is independent of the choice of $s_i$: the product with its
conjugate is the nonzero square $B_i z_i^2$ modulo $p$.
For a raw word $a$, its central Frobenius sign is
$\prod_i\varepsilon_i(p)^{a_i}$. The prime splits completely in $N$
exactly when every word in \eqref{ar:completed-words} has sign $+1$.
The analogous test for $M$ uses \eqref{ar:full-words}. Any detected
central sign gives residue degree two in that field. Since $N\subset M$,
the completed-field detections are a subset of the full-field detections.
Order the seven completed predicates before the twelve full predicates,
and assign each prime to its first nonzero predicate. This gives
disjoint sets of primes detected in the two fields, which we will use
for the finite prime corrections in the analytic estimate.

\section{The Relative Zeta Factor and Its Explicit Ceiling}
\label{an:section}

We bound the relative zeta value needed in the geometric construction.
Each zeta function is normalized by the absolute degree of its field.
With this normalization, the relative degree contributes no additional
factor of two to the final bound.

Let $K$ be a field in Theorem~\ref{tw:field-family}, let $c_0$ be its
actual complex conjugation, and put
\[
 F=K^{\langle c_0\rangle},\qquad [K:\mathbb Q]=2d,\qquad
 [F:\mathbb Q]=d=b+2c,\qquad \theta=c/d.
\]
The extension $K/F$ is quadratic and unramified at all finite primes.
The field $K$ is totally imaginary and $b\geq1$. Its root discriminant, equal to
that of $F$, is
\[
 \lambda=2^{9/4}\sqrt{15015},\qquad
 \ell=\log\lambda=\frac94\log2+\frac12\log15015.
\]
Write $L_F(s)=\zeta_K(s)/\zeta_F(s)$. We will prove the following bound.

\begin{theorem}\label{an:ceiling}
For all sufficiently large degrees in the constructed field family,
\[
 \frac{\log L_F(1)}d<0.042161819.
\]
No generalized Riemann hypothesis, effective lower bound for the first
such degree, or asymptotic class-number formula is required.
\end{theorem}

\subsection{Finite-Field Comparison and Passage to One}

For a number field $A$, denote its class number, regulator, and number
of roots of unity by $h_A,R_A,w_A$. The analytic class-number formula
gives the positive value
\begin{equation}\label{an:class-formula}
 L_F(1)=2^c\pi^{d-c}\lambda^{-d/2}
       \frac{h_KR_Kw_F}{h_FR_Fw_K}>0.
\end{equation}
Indeed, the residues of $\zeta_A$ are
$2^{r_1(A)}(2\pi)^{r_2(A)}h_AR_A/(w_A\sqrt{|\Delta_A|})$,
and finite unramifiedness gives $|\Delta_K|=|\Delta_F|^2$.
We bound the factor in~\eqref{an:class-formula}. The relative unit and
ideal-class normalizations enter later in the geometric argument.

At a prime of $F$ of norm $q$, quadratic splitting gives respectively
\[
 \zeta_{K,v}(s)=(1-q^{-s})^{-2}
 \quad\hbox{or}\quad
 \zeta_{K,v}(s)=(1-q^{-2s})^{-1}.
\]
Both are at most $\zeta_{F,v}(s)^2$ for real $s>1$. Consequently
\begin{equation}\label{an:relative-euler}
 \frac{\log L_F(s)}d
 \leq\frac{\log\zeta_K(s)}{2d}.
\end{equation}
For a rational prime $p$ with absolute ramification index $e$ and
residue degree $f$ in a Galois field, its contribution to the normalized
logarithm is
\begin{equation}\label{an:local-phi}
 \Phi_{p,s}(e,f)=\frac{-\log(1-p^{-fs})}{ef}.
\end{equation}
Multiplying either $e$ or $f$ by a positive integer cannot increase
this expression. Hence, in an extension of finite Galois fields, the
normalized zeta function of the larger field is at most that of the
smaller field, prime by prime.

Set $\Gamma_{\mathbb R}(s)=\pi^{-s/2}\Gamma(s/2)$ and
$\Gamma_{\mathbb C}(s)=2(2\pi)^{-s}\Gamma(s)$. The nontrivial
quadratic Hecke character of $K/F$ is unramified at finite places and
odd at each real place of $F$. The Hecke continuation and functional
equation theorem \cite{Hecke1920,Tate1967} gives, up to a
positive constant independent of $s$, the completion
\begin{equation}\label{an:completion}
 \mathcal L_F(s)=\lambda^{ds/2}\pi^{-bs/2}(2\pi)^{-cs}
       \Gamma((s+1)/2)^b\Gamma(s)^cL_F(s).
\end{equation}
It is entire of order one. Taking the quotient of the completed
Dedekind zeta functional equations shows that its root number is $+1$.

The zero pairing in the functional equation lets us compare this
completion at one with its values to the right of one.
\begin{lemma}\label{an:monotonicity}
The function $\mathcal L_F(s)$ is nondecreasing on the real interval
$[1,\infty)$.
\end{lemma}
\begin{proof}
The Euler product and functional equation place every zero in
$0\leq\operatorname{Re}s\leq1$. Center the function at $1/2$ and
write $t=s-1/2$. Its zeros occur in pairs $z,-z$ and in conjugate
pairs. Since its order is one, $\sum_{z\ne0}|z|^{-2}<\infty$.
It follows that the paired canonical product and its logarithmic
derivative converge locally off the zeros. Evenness makes the remaining linear exponential
factor constant. A zero of multiplicity $m$ at the center contributes
$m/t$ to the logarithmic derivative. A pair with $z=a+iv$ contributes
the real part
\[
 \operatorname{Re}\frac{2t}{t^2-z^2}
 =\frac{2t(t^2-a^2+v^2)}
 {(t^2-a^2+v^2)^2+4a^2v^2}\geq0\qquad(t\geq1/2),
\]
because $|a|\leq1/2$. There are no zeros on $[1,\infty)$:
this follows from the Euler product above one and
\eqref{an:class-formula} at one. The completion is positive there.
Its logarithmic derivative is nonnegative, proving the lemma.
\end{proof}

For $\varepsilon>0$, division of the two completion factors gives
\begin{align}
 \frac{\log L_F(1)}d
 &\leq\frac{\log\zeta_K(1+\varepsilon)}{[K:\mathbb Q]}
           +G(\varepsilon,\theta),\label{an:finite-gamma}\\
 G(\varepsilon,\theta)
 &=\frac\varepsilon2(\ell-\log\pi)-\varepsilon\theta\log2
 +(1-2\theta)\log\Gamma(1+\varepsilon/2)
 +\theta\log\Gamma(1+\varepsilon).
 \label{an:gamma}
\end{align}
In particular $G(\varepsilon,\theta)\to0$ uniformly for
$0\leq\theta\leq1/2$ as $\varepsilon\downarrow0$.

\subsection{An Unconditional Asymptotic Prime Budget}

The basic inequality for limiting prime counts improves a finite Euler
estimate at a fixed abscissa to the following bound at one. Only this
inequality is needed, without a Brauer--Siegel equality.

\begin{proposition}\label{an:prime-budget}
Suppose that, at $\sigma=1+\varepsilon>1$, every field in the family
satisfies
\[
 \frac{\log\zeta_K(\sigma)}{[K:\mathbb Q]}\leq Y_K(\sigma)
 \leq Y_*(\sigma).
\]
Let $e_p^0,f_p^0$ be uniform lower bounds for its absolute local indices,
and suppose
\[
 R\geq\sum_p\frac{\log p}{e_p^0(p^{2f_p^0}-1)}.
\]
Then
\begin{equation}\label{an:budget-conclusion}
 \limsup_{d\to\infty}\frac{\log L_F(1)}d
 \leq Y_*(1+\varepsilon)
 +\varepsilon\left(\frac{\ell-\gamma_E-\log(4\pi)}4+R\right).
\end{equation}
Every strictly larger constant is an eventual uniform ceiling.
\end{proposition}
\begin{proof}
Take any sequence with $n=[K:\mathbb Q]=2d\to\infty$. A diagonal
subsequence makes each normalized prime count converge:
$\beta_q=\lim N_q(K)/n$, where $q$ ranges over prime powers.
The fixed root discriminant makes this an asymptotically exact family.
The unconditional Basic Inequality$'$ of
\cite[Proposition 3.1]{TsfasmanVladut2002} is
\begin{equation}\label{an:tv-budget}
 \sum_q\beta_q w(q)\leq B,\qquad
 w(q)=2\log q\sum_{m\geq1}\frac1{q^m+1},\qquad
 B=\frac{\ell-\gamma_E-\log(4\pi)}2.
\end{equation}
To check the normalization, their genus is $g=n\ell/2$, their finite
invariants are $\phi_q=2\beta_q/\ell$, and total imaginary signature
gives $\phi_{\mathbb R}=0$, $\phi_{\mathbb C}=1/\ell$.
The proposition is unconditional, with archimedean constant
$\gamma_E+\log(4\pi)$.

Put $a_s(q)=-\log(1-q^{-s})$ and
$Z_\infty(s)=\sum_q\beta_q a_s(q)$. The series converges at one,
since $a_1(q)/w(q)\leq(q+1)/(2(q-1)\log q)$ is bounded for $q\geq2$.
For every fixed $\tau>1$, dominated convergence in the normalized
Euler products gives
\[
 \lim\frac{\log\zeta_K(\tau)}n=Z_\infty(\tau).
\]
One can dominate prime by prime by $-\log(1-p^{-\tau})$.
Equation~\eqref{an:finite-gamma}, followed by $\tau\downarrow1$,
then gives $\limsup d^{-1}\log L_F(1)\leq Z_\infty(1)$.
The uniform vanishing of $G$ and monotone convergence justify this
order of limits. The finite-field Euler products are used only to the
right of one.

For every prime power $q$,
\begin{align*}
 a_1(q)-a_{1+\varepsilon}(q)&\leq
                \frac{\varepsilon\log q}{q-1},\\
 \frac{\log q}{q-1}-\frac{w(q)}2
 &=\log q\sum_{m\geq1}\frac1{q^m(q^m+1)}
 \leq\frac{\log q}{q^2-1}.
\end{align*}
For a Galois field, the normalized contribution above $p$ to the last
weighted sum is $\log p/[e_p(p^{2f_p}-1)]$. The local floors and
Fatou's lemma bound its limiting sum by $R$. Combining this with
\eqref{an:tv-budget} proves
$Z_\infty(1)\leq Z_\infty(1+\varepsilon)+\varepsilon(B/2+R)$.
This proves the assertion on every such subsequence. If a strictly
larger constant were not an eventual uniform bound, a sequence of
violating fields would have a diagonal subsequence of the kind just
considered, giving a contradiction.
\end{proof}

\subsection{The Complete Finite-Field Euler Factorization}
\label{an:factorization}

Let $N$ be the degree-$16384$ retained field of
Proposition~\ref{ar:central-fields}, and set
$S'=\{2,3,5,7,11,13,17\}$. A superscript $S'$ removes these Euler
factors. Proposition~\ref{ar:hecke-data} gives all irreducible
representations of $\operatorname{Gal}(N/\mathbb Q)$ and realizes
each nonlinear factor as an induced nontrivial finite-order Hecke
character. There are $128$ linear, $256$ quadratic, $456$ quartic,
and $124$ octic factors. Among the quartics, eight have pure gamma
parity, and every other nonlinear factor is mixed. The regular
representation identity is
\begin{equation}\label{an:regular-factorization}
 \zeta_N(s)=\prod_\rho L(s,\rho)^{\dim\rho},\qquad
 128+256\cdot4+456\cdot16+124\cdot64=16384.
\end{equation}
For real $s>1$ the factors are nonzero. Real factors are positive:
their Euler products tend to one as $s\to\infty$ and cannot cross
zero. Conjugate factors have product $|L(s,\rho)|^2$. Thus
\begin{equation}\label{an:Y}
 Y(s):=\frac{\log\zeta_N^{S'}(s)}{16384}
 =\frac1{16384}\sum_\rho\dim\rho\,
          \log|L^{S'}(s,\rho)|.
\end{equation}
Each conjugate pair is counted twice in this sum, with the same
representation degree. An upper bound for each modulus may be
substituted after this identity, even when some logarithms are negative.

The exact arithmetic input for one factor consists of its conductor
$Q$, gamma parities, root number when proved to be one, an inducing
radicand, genus twist, and the exceptional denominators
\[
 D_{\rho,p}(T)=\det(1-T\rho(\operatorname{Frob}_p)
                         \mid V_\rho^{I_p}).
\]
Their construction, including the cancellation of quartic twists at
5 and 13, is given in Proposition~\ref{ar:hecke-data}. The
analytic evaluations below use the representations constructed there,
including their conductors and local characters. Formal compatibility
among a set of Euler polynomials would not identify these functions.

\subsection{Exact Coefficients and Their Enumeration}
\label{an:coefficients}

Write $L(s,\rho)=\sum_{n\geq1}a_n n^{-s}$ and $r=\dim\rho$.
At every prime, the reciprocal roots of $D_{\rho,p}$ have modulus
one, and there are at most $r$ of them. Hence
\begin{equation}\label{an:divisor-bound}
 |a_n|\leq d_r(n),\qquad
 \sum_{n\geq1}d_r(n)n^{-\beta}=\zeta(\beta)^r
 \quad(\beta>1).
\end{equation}
The finite-image representation also makes every coefficient used here
an exact Gaussian integer.

To construct the coefficients outside $S'$, let
$v_p\in\mathbb F_2^7$ be genus Frobenius, and let $B_q$ be the polar
form of the corresponding central character. Outside its radical,
Frobenius has eigenvalues in opposite pairs and square $\kappa I$,
where $\kappa\in\{1,-1\}$. Its denominator is
$(1-\kappa T^2)^{r/2}$, before the explicitly supplied quartic twist.
The coefficients at exponents $2j$ are
$\binom{j+r/2-1}{r/2-1}\kappa^j$, and odd exponents vanish.
Inside the radical, Frobenius is scalar, and its denominator is
$(1-\omega T)^r$, where $\omega\in\{1,-1,i,-i\}$ is determined
by the inducing field and twist. Its coefficients are
$\binom{j+r-1}{r-1}\omega^j$.

At a scalar prime the inducing multiquadratic base splits completely.
For the positive quadratic induction, the scalar sign is obtained by
evaluating the radicand at these embeddings modulo $p$ and taking
quadratic residue symbols. These symbols agree
because the represented Frobenius is scalar. There is an exact
procedure even at a prime dividing the norm of the displayed radicand.
For an integral radicand $\eta$, let $m=v_p(N\eta)$. Each embedding
valuation is between zero and $m$. Lift the base square roots to
modulus $p^{m+1}$ by Hensel's lemma. This determines every embedding
valuation and normalized nonzero unit residue. The valuations are
even, by unramifiedness outside $S'$. The normalized quadratic symbols
give the required signs. Rational square normalization does not change
the extension. This procedure treats every norm prime, rather than
assuming that a finite list of ordinary exceptions is exhaustive.

The genus and twist characters are evaluated exactly by Kronecker
symbols and by discrete logarithms in $(\mathbb Z/5\mathbb Z)^\times$
or $(\mathbb Z/13\mathbb Z)^\times$, with $\xi_p(2)=i$.
The genus characters have common period $120120$, so lookup tables may
use this period or any multiple of it. At $p\in S'$, the
verified polynomial $D_{\rho,p}$ supplies the recurrence
\[
 a_{p^j}=-\sum_{k=1}^{\min(j,r)}[T^k]D_{\rho,p}(T)\,a_{p^{j-k}},
 \qquad a_1=1.
\]
Multiplicativity then determines all coefficients.

To enumerate the nonzero coefficients through the cutoff $N_0$, first
list all primes up to $N_0$ by a segmented sieve using every sieving
prime through $\lfloor\sqrt{N_0}\rfloor$. Assign an ordinary nonscalar prime the
first possible support location $p^2$, a scalar prime the location $p$,
and an exceptional prime the safe lower bound $p$. Sort primes by these
locations. Recursion appends only later primes in that order, considering
every admissible prime power and continuing across intermediate zero
coefficients. Unique factorization then visits each supported integer
exactly once. If a first support location exceeds the remaining quotient,
that prime and every later prime can be skipped. These observations prove that
the enumeration is complete.

Signed and absolute moments use arbitrary-precision integers. Fixed-width
coefficient arithmetic is checked before conversion, and intermediate
products are formed in a wider type. An overflow aborts the calculation.
Modular multiplication uses a wide product or a checked Montgomery
reduction as appropriate. The coefficients, support counts, and moments are computed exactly.

\subsection{The Linear Factors}

The $128$ linear factors are the primitive quadratic Dirichlet functions
of the genus field
$E=\mathbb Q(i,\sqrt2,\sqrt3,\sqrt5,\sqrt7,\sqrt{11},\sqrt{13})$,
including the principal factor $\zeta(s)$. For a signed squarefree
product $D_0$ of these radicands, its fundamental discriminant is
$D=D_0$ if $D_0\equiv1\pmod4$, and $D=4D_0$ otherwise.
The character is the Kronecker symbol $(D/\cdot)$.
The implementation checks its primitive conductor and its values on
explicit generators of every prime-power component of the unit group.
These values identify the character independently of its library index.

The finite identity
\begin{equation}\label{an:dirichlet-hurwitz}
 L(s,\chi)=q^{-s}\sum_{a=1}^{q}\chi(a)\zeta(s,a/q)
\end{equation}
reduces evaluation to Hurwitz zeta functions. Euler--Maclaurin summation
gives the explicit enclosure
\begin{align*}
 \zeta(s,x)={}&\sum_{k=0}^{M-1}(k+x)^{-s}
 +\frac{(M+x)^{1-s}}{s-1}+\frac12(M+x)^{-s}\\
 &+\sum_{j=1}^{J}\frac{B_{2j}}{(2j)!}(s)_{2j-1}
                     (M+x)^{1-s-2j}+\mathcal R,\\
 |\mathcal R|\leq{}&
 \frac{2\zeta(2J)(s)_{2J}}{(2\pi)^{2J}(s+2J-1)}
                       (M+x)^{1-s-2J},
 \qquad s>1, x>0.
\end{align*}
The remainder follows from the periodic Bernoulli bound
$|\widetilde B_{2J}(t)|\leq2(2J)!\zeta(2J)/(2\pi)^{2J}$.
For rigorous high-precision evaluation by Euler--Maclaurin summation,
see \cite{Johansson2015Hurwitz}. The supplied evaluation uses Arb/FLINT
ball arithmetic \cite{Johansson2017}, retaining the endpoints of each
real value.
Removing $S'$ multiplies a factor by
$\prod_{p\in S'}(1-\chi(p)p^{-s})$.

\subsection{An Approximate Functional Equation for Mixed Factors}
\label{an:mixed-afe}

Numerical evaluation from gamma factors and a functional equation is
developed in \cite{Dokchitser2004}. We use this approach with explicit
kernel and tail bounds for the factors constructed above.
For a mixed factor of degree $2m$, where $m=1,2,4$, its completion is
\[
 \Lambda(s)=Q^{s/2}\Gamma_{\mathbb R}(s)^m
                    \Gamma_{\mathbb R}(s+1)^m L(s)
 =2^m Q^{s/2}(2\pi)^{-ms}\Gamma(s)^mL(s).
\]
It is entire and satisfies
$\Lambda(s)=\varepsilon_\rho\overline{\Lambda(1-\overline s)}$,
with $|\varepsilon_\rho|=1$. For a real quadratic induction the
Dedekind zeta quotient gives $\varepsilon_\rho=1$.

Let $k_m(x)$ have Mellin transform $\Gamma(z)^m$, and put
\[
 H_{m,a}(x)=\int_1^\infty u^{a-1}k_m(xu)\,du,\qquad
 t=\frac{(2\pi)^m}{\sqrt Q},\qquad
 F_m=\frac{t^s}{\Gamma(s)^m}.
\]
These kernels are positive. More explicitly,
\begin{equation}\label{an:positive-kernel}
 k_m(x)=\int_{(0,\infty)^{m-1}}
 \exp\left(-u_1-\cdots-u_{m-1}-\frac{x}{u_1\cdots u_{m-1}}\right)
                  \frac{du_1\cdots du_{m-1}}{u_1\cdots u_{m-1}},
\end{equation}
with $k_1(x)=e^{-x}$. Taking the Mellin transform under this positive
integral proves the asserted transform and also complete monotonicity
in $x$. In particular $k_2(x)=2K_0(2\sqrt x)$
\cite[10.32.10, 10.43.19]{DLMF}.

Splitting the Mellin integral at one gives a common approximate
functional equation for the mixed factors.
\begin{lemma}\label{an:balanced-afe}
For real $s>1$ and every nonlinear mixed factor,
\begin{equation}\label{an:mixed-formula}
 L(s)=F_m\sum_{n\geq1}a_nH_{m,s}(tn)
 +\varepsilon_\rho F_m\sum_{n\geq1}\overline{a_n}
                                      H_{m,1-s}(tn).
\end{equation}
\end{lemma}
\begin{proof}
The function $\Theta(u)=\sum a_n k_m(tnu)$ has Mellin transform
$\Gamma(z)^m t^{-z}L(z)=2^{-m}\Lambda(z)$. Mellin inversion and
the completed functional equation give
$\Theta(u)=\varepsilon_\rho u^{-1}\overline{\Theta(1/u)}$.
The contour shift has no residues because the completion is entire,
and the gamma factors give the needed vertical decay. Splitting the Mellin
integral at one and substituting $u\mapsto1/u$ below one yields
\eqref{an:mixed-formula}. The positive integral representation and
\eqref{an:divisor-bound} justify the sums in either exponentially
decaying half.
\end{proof}

For degree two this uses
$H_{1,a}(x)=\int_1^\infty u^{a-1}e^{-xu}\,du$.
For the mixed quartic implementation it is convenient to write
\[
 I_b(y)=\int_1^\infty u^bK_0(yu)\,du,\qquad
 H_{2,a}(x)=4I_{2a-1}(2\sqrt x).
\]
Its prefactor is $4t^s/\Gamma(s)^2$ when the stored kernel
is $I_{2s-1}+I_{1-2s}$. This identity accounts for the factor four
in the quartic evaluation. For octics the stored kernel is
$H_{4,a}$ itself, with prefactor $t^s/\Gamma(s)^4$.

\subsection{Tail Bounds and Absolute Coefficient Prefixes}
\label{an:tails}

For $\beta\geq\max(s,1-s)$,
$H_{m,a}(x)\leq H_{m,\beta}(x)$ when $a\leq\beta$, and
\[
 x^\beta H_{m,\beta}(x)=\int_x^\infty v^{\beta-1}k_m(v)\,dv
\]
is decreasing. If $M_{N_0}(\beta)$ bounds
$\sum_{n>N_0}|a_n|n^{-\beta}$, either omitted half of
\eqref{an:mixed-formula} is at most
\begin{equation}\label{an:general-tail}
 F_m (N_0+1)^\beta H_{m,\beta}(t(N_0+1))M_{N_0}(\beta).
\end{equation}
One may always use $M_{N_0}(\beta)=\zeta(\beta)^{2m}$.

The quadratic and quartic evaluations also use elementary envelopes.
For the exponents in question, with $1<s<1.01$,
\begin{align}
 H_{1,s}(x),H_{1,1-s}(x)&\leq
 V(x):=e^{-x}(1+1/x)/x,\label{an:exp-envelope}\\
 I_{2s-1}(y),I_{1-2s}(y)&\leq
 U(y):=\sqrt{\pi/2}\,e^{-y}(1+1/y)y^{-3/2}.
 \label{an:bessel-envelope}
\end{align}
The first follows by replacing $u^{s-1}$ and $u^{-s}$ by $u$.
For the second use $K_0(yu)\leq\sqrt{\pi/(2yu)}e^{-yu}$ and
$u^{b-1/2}\leq u$ for $b\leq3/2$.
The functions $n^\beta V(tn)$ and
$n^\beta U(2\sqrt{tn})$ decrease respectively when
$tn>\beta-1$ and $2\sqrt{tn}>2\beta-3/2$.
Thus substituting these envelopes and the divisor majorant gives
valid tails even when the chosen $\beta$ exceeds an exponent covered
by the envelope. The quartic prefactor includes the factor four from the conversion
between the kernels. When we sum both halves, we include both tail
allowances.

For the octics a sharper absolute series is available. Let
$W=\operatorname{rad}B_q$. Then $[E^W:\mathbb Q]=64$. At an ordinary
prime, absolute coefficients have local series $(1-T)^{-8}$ if
$v_p\in W$, and $(1-T^2)^{-4}$ otherwise. These are precisely the
local factors of $\zeta_{E^W}^{S'}(\beta)^{1/8}$.
For every exceptional denominator in the supplied arithmetic table,
an exact Gaussian-polynomial comparison supplies nonnegative integers
$u_p,v_p$ and a fourth root of unity $\omega_p$ such that
\begin{equation}\label{an:exceptional-phase}
 D_{\rho,p}(T)=(1-\omega_pT)^{u_p}
                          (1-\omega_p^2T^2)^{v_p}.
\end{equation}
The check enumerates $u_p,v_p\geq0$ with $u_p+2v_p\leq8$ and
the four possible $\omega_p$, comparing all polynomial coefficients
exactly. After the rotation $T\mapsto\omega_p^{-1}T$, the reciprocal
has nonnegative coefficients. Consequently the full absolute series is
exactly
\begin{equation}\label{an:absolute-series}
 M_\rho(\beta):=\sum_{n\geq1}|a_n|n^{-\beta}
 =\zeta_{E^W}^{S'}(\beta)^{1/8}
 \prod_{p\in S'}(1-p^{-\beta})^{-u_p}
                       (1-p^{-2\beta})^{-v_p}.
\end{equation}
The same rotations show that each nonzero coefficient lies on a
Gaussian coordinate axis, so $|a_n|=|\operatorname{Re}a_n|+
|\operatorname{Im}a_n|$ exactly. The genus subfield zeta value is
evaluated using its $64$ quadratic characters by
\eqref{an:dirichlet-hurwitz}.

We subtract a rigorous lower bound for the finite absolute prefix.
On a bin with integer midpoint $m_0$ and radius $r_0$, record exact
moments $A_j=\sum|a_n|(n-m_0)^j$, $0\leq j\leq J=20$.
Let $h=\lceil\beta\rceil$ and $\rho_0=r_0/m_0<1/16$.
The degree-$J$ binomial expansion of $n^{-\beta}$ has error at most
\begin{equation}\label{an:absolute-prefix-error}
 A_0m_0^{-\beta}\binom{J+h}{h-1}
                    \frac{\rho_0^{J+1}}{(1-\rho_0)^h}.
\end{equation}
Indeed $(\beta)_j/j!\leq\binom{j+h-1}{h-1}$, and after the first
omitted coefficient their ratios are bounded by the coefficients of
$(1-\rho_0)^{-h}$. Sum the interval polynomials and subtract these
errors to obtain $P_{N_0}(\beta)\leq\sum_{n\leq N_0}|a_n|n^{-\beta}$.
Then $M_\rho(\beta)-P_{N_0}(\beta)$ is an admissible
$M_{N_0}(\beta)$ in~\eqref{an:general-tail}.
The octic evaluation selects the smallest valid enclosure among
$\beta=5/4,3/2,2,5/2,3$. Each candidate is an upper bound, so this selection preserves the
inequality.

\subsection{The Eight Pure Quartic Factors}
\label{an:pure}

The pure factors in Proposition~\ref{ar:hecke-data} are real quadratic
Hecke inductions with root number one, conductor $Q=240240^2$, and
gamma vector $(a,a,a,a)$, $a\in\{0,1\}$, four of each parity.
Their completion is $Q^{s/2}\Gamma_{\mathbb R}(s+a)^4L(s)$.
Put
\[
 t=\pi^2/\sqrt Q,\quad b_0=(s+a)/2,\quad b_1=(1-s+a)/2,
 \quad F=\frac{t^{s+a}}{\Gamma(b_0)^4},\quad K_b=H_{4,b}.
\]
Then the exact functional equation is
\begin{equation}\label{an:pure-afe}
 L(s)=F\sum_{n\geq1}a_n n^a
       \bigl(K_{b_0}(t^2n^2)+K_{b_1}(t^2n^2)\bigr).
\end{equation}
To verify the normalization, the Mellin transform of
$\sum a_n n^a k_4(t^2n^2u)$ is
\[
 \Gamma(z)^4t^{-2z}L(2z-a)=Q^{a/2}\Lambda(2z-a).
\]
Reflection is about $z=(a+1/2)/2$. Splitting at one gives
\eqref{an:pure-afe}, with no extra factor two.
For $\beta\geq\max(s,1-s)$, put $B=(\beta+a)/2$. Either omitted
side is bounded by
\begin{equation}\label{an:pure-tail}
 F(N_0+1)^{\beta+a}K_B(t^2(N_0+1)^2)\zeta(\beta)^4.
\end{equation}
This follows from the same decreasing-integral argument as
\eqref{an:general-tail}. The exact cutoff here is
$N_0=8\sqrt Q=1921920$.

\subsection{Kernel and Moment Evaluation}
\label{an:kernel-enclosures}

We evaluate the kernels on finite grids and bound the interpolation
errors. For the four-gamma kernel, let
$H_n^{(j)}=\sum_{k=1}^n k^{-j}$ and
\[
 L_n=4(H_n^{(1)}-\gamma_E)-\log x,\quad
 P_{3,n}=L_n^3/6+2L_n(\zeta(2)+H_n^{(2)})
                     +4(H_n^{(3)}-\zeta(3))/3,
\]
\[
 P_{2,n}=L_n^2/2+2(\zeta(2)+H_n^{(2)}).
\]
Residues at the order-four poles of $\Gamma(z)^4$ give the convergent
series
\begin{align}
 k_4(x)&=\sum_{n\geq0}\frac{x^n}{(n!)^4}P_{3,n},
 \label{an:k4-series}\\
 K_a(x)&=\Gamma(a)^4x^{-a}-
 \sum_{n\geq0}\frac{x^n}{(n!)^4}
 \left(\frac{P_{3,n}}{n+a}+\frac{P_{2,n}}{(n+a)^2}
       +\frac{L_n}{(n+a)^3}+\frac1{(n+a)^4}\right).
 \label{an:K4-series}
\end{align}
The second identity initially follows for $a>0$ and extends to
$-1/100<a<0$ by analytic continuation. None of the exponents evaluated
here is zero or a negative integer, so no colliding-pole formula is
needed.

To bound the tails of these series, for $n\geq2$ set
\[
 B_n=|\log x|+4(1+\log n+\gamma_E),\quad
 U_{3,n}=B_n^3/6+4\zeta(2)B_n+4\zeta(3)/3,
 \quad U_{2,n}=B_n^2/2+4\zeta(2).
\]
They bound $|L_n|,|P_{3,n}|,|P_{2,n}|$ in the corresponding order.
Also $B_{n+1}\leq(1+1/n)B_n$, so
\[
 r_n=\frac{x}{(n+1)^4}(1+1/n)^3
\]
is an upper bound for the subsequent term ratios of the majorants.
Once $r_n<1$, the entire omitted tail of~\eqref{an:k4-series}
is at most
\[
 \frac{x^n}{(n!)^4}\frac{U_{3,n}}{1-r_n}.
\]
For~\eqref{an:K4-series}, provided $a>-1/100$, it is at most
\[
 \frac{x^n}{(n!)^4}\frac1{1-r_n}
 \left(\frac{U_{3,n}}{n-1}+\frac{U_{2,n}}{(n-1)^2}
       +\frac{B_n}{(n-1)^3}+\frac1{(n-1)^4}\right).
\]
All these expressions are evaluated outward. The first four Taylor
coefficients of $k_4$ are obtained by the same argument on
$|z-x|\leq x/2$, replacing $x$ by $3x/2$ in the numerator bound
and $|\log x|$ by $|\log x|+\log2+\pi/2$.
Cauchy's estimate bounds the error in the $j^{th}$ coefficient by
the uniform tail times $(2/x)^j$.

Higher Taylor coefficients follow from
\begin{equation}\label{an:kernel-odes}
 (x\partial_x)^4k_4=xk_4,\qquad xK_a'+aK_a=-k_4.
\end{equation}
These identities follow directly from the Mellin transform and the
defining integral. They give finite recurrences using rational
operations and the enclosed initial values. The two-gamma kernel
uses $K(x)=K_0(2\sqrt x)$, $J_b(x)=I_b(2\sqrt x)$, with
\[
 xK''+K'-K=0,\qquad 2xJ_b'+(b+1)J_b=-K.
\]
If $k_j,j_j$ denote Taylor coefficients at $x_0$, the recurrences are
\[
 k_{j+2}=\frac{k_j-(j+1)^2k_{j+1}}{x_0(j+1)(j+2)},\qquad
 j_{j+1}=\frac{-k_j-(2j+b+1)j_j}{2x_0(j+1)}.
\]
The initial values are $K_0(2\sqrt{x_0})$ and
$-K_1(2\sqrt{x_0})/\sqrt{x_0}$. The one-gamma case uses the elementary
exponential and $xH_{1,a}'+aH_{1,a}=-e^{-x}$.

The grids have consecutive endpoints $x_0,31x_0/32$.
For any positive Laplace kernel $H$ in use, its holomorphic extension
obeys $|H(z)|\leq H(x_0/2)$ on $|z-x_0|\leq x_0/2$.
Hence a degree-$J$ Taylor polynomial has error at most
\begin{equation}\label{an:taylor-error}
 H(x_0/2)\frac{16^{-J-1}}{1-1/16}
\end{equation}
throughout the bin. The backward Taylor remainder is nonnegative by
complete monotonicity, allowing propagation of the exponential and
Bessel kernels from a certified large-argument interval.
Their starting values are bounded by~\eqref{an:exp-envelope} and
\eqref{an:bessel-envelope}. The four-gamma kernel instead uses the
residue series at each center. Grid endpoints, their integer floors,
and complete coverage from $1$ to the exact coefficient cutoff are
checked with outward intervals. The calculation is rejected if an
integer floor is unresolved.

For a linear kernel argument $x=tn$, translating its degree-$20$
Taylor polynomial at $x_0$ to the integer midpoint $m_0$ gives
\[
 [ (n-m_0)^j ]\,P(tn)
 =t^j\sum_{h=j}^{20}[ (x-x_0)^h ]P(x)
                 \binom hj(tm_0-x_0)^{h-j}.
\]
Contract this polynomial against the exact real and imaginary moments.
Multiplying~\eqref{an:taylor-error} by the sum of absolute coefficients
in the bin bounds its contribution to the error.
For the pure quartics the substitution is quadratic:
\[
 t^2n^2-x_0=(t^2m_0^2-x_0)+2t^2m_0(n-m_0)+t^2(n-m_0)^2.
\]
A degree-$20$ kernel polynomial thus needs moments through degree
$40$. These moments and their absolute masses include $n^a$ exactly
once. The uniform remainder still has degree $20$ in the kernel
variable, so~\eqref{an:taylor-error} applies before composition.

Suppose $A,B$ are the finite primary and conjugated dual sums, and
$E_A,E_B$ bound their interpolation and infinite-tail errors. A known
root number one gives $|L(s)|\leq|A+B|+E_A+E_B$. With an unknown
phase one always has $|L(s)|\leq|A|+|B|+E_A+E_B$.
Every modulus, gamma value, logarithm, and conversion to an endpoint
retains its rounding enclosure.

\subsection{Recovering the Mixed Quadratic and Quartic Phase}
\label{an:two-cutoffs}

For the complex quadratic and mixed quartic factors, a second Mellin
cutoff can sharpen this bound by enclosing the root number. Split the
Mellin integral at $y>0$ instead of one. The two exact terms are
\[
 A_y=F_m y^s\sum a_nH_{m,s}(tyn),\qquad
 B_y=F_m y^{s-1}\sum\overline{a_n}H_{m,1-s}(tn/y),
 \qquad L(s)=A_y+\varepsilon_\rho B_y.
\]
Take $y=(31/32)^{22}$, so shifted kernels reuse the same exact moment
bins. Any terminal part outside the stored kernel range is put into
the tail with its actual first omitted integer.

For finite approximations, write $U=A_y-A_1$, $V=B_1-B_y$ and let
$\eta,\Delta$ bound their respective errors. If $|V|>\Delta$, then
\begin{equation}\label{an:phase-disk}
 \varphi=U/V,\qquad
 |\varepsilon_\rho-\varphi|
 \leq r_\varphi:=\frac{\eta+|\varphi|\Delta}{|V|-\Delta}.
\end{equation}
This follows by writing the exact quotient as
$(U+u)/(V+v)$, with $|u|\leq\eta$, $|v|\leq\Delta$.
If each balanced side has error at most $E_1$, the resulting bound is
\begin{equation}\label{an:phase-upper}
 |L(s)|\leq|A_1+\varphi B_1|+2E_1+|B_1|r_\varphi.
\end{equation}
The coefficient $2E_1$ uses the known identity
$|\varepsilon_\rho|=1$, not an assumption that $|\varphi|=1$.
All distances from zero and radii are enclosed outward. If the
denominator cannot be separated from zero, the phase-free bound is
retained. Complex octic factors use the phase-free bound directly.

\subsection{The Certified Factor Sums}
\label{an:factor-sums}

The arithmetic table and exact moment files in the supplementary
certificate determine every factor just described. The abscissa is
the exact rational
\[
 \sigma=\frac{12001}{12000}.
\]
For each row, the evaluator applies the appropriate AFE, includes
interpolation and both infinite tails, and removes the exceptional
factors by
\[
 \log|L^{S'}(\sigma,\rho)|
 \leq\log U_\rho+
       \sum_{p\in S'}\log|D_{\rho,p}(p^{-\sigma})|,
\]
where $U_\rho$ is its proved modulus upper bound.
The octic cutoffs are
$N_0=\min(3\cdot10^9,\lfloor\sqrt Q/40\rfloor)$.
The exact integers for the quadratic and mixed quartic rows are
specified in their moment records and individually checked when
constructing the bins and tails. The pure quartic cutoff was given
in Section~\ref{an:pure}.

The following table gives rational upper bounds for the sums of
logarithms. The number of factors includes both members of each
conjugate pair. An interval enclosing an AFE upper bound encloses that
bound, and need not give a two-sided enclosure of comparable width for
the $L$-value itself.
\begin{center}
\begin{tabular}{lrr}
\toprule
Type & Number & Upper bound for $\sum\log|L^{S'}(\sigma,\rho)|$\\
\midrule
Linear &128& $7.337578130846530221$\\
Mixed quadratic &256& $0.400886873183882259$\\
Mixed quartic &448& $-0.046453348154273950$\\
Pure quartic &8& $0.058145052434373409$\\
Mixed octic &124& $-0.111176963043772737$\\
\bottomrule
\end{tabular}
\end{center}
Exact addition of their dimension-weighted upper endpoints in
\eqref{an:Y} gives
\begin{equation}\label{an:Y-numerical}
 Y(\sigma)\leq0.00044535540710354679437255859375
             <0.000445355407103547.
\end{equation}
The coefficient algorithms, kernel bounds, and tail estimates above
specify how to verify these inequalities. The supplementary program
checks the exact row identities and partitions before performing the
interval evaluations. Its moment mode regenerates the newest 64 sector
moment files from the coefficient algorithms. Earlier moment records
remain finite inputs, and the data manifest checks their identities.
% CHECK: The current replay does not regenerate all earlier moments or actual-field data.
% The required independent verification scope is recorded in evidence/body-review-arithmetic-20260927.md.

\subsection{Finite Corrections beyond the Retained Field}
\label{an:census}

To compare $N$ with the whole tower, restore $S'$ using the actual
local indices of $K$. Its exact contribution is
\begin{align}
 B_D(\sigma)={}&-\frac{\log(1-2^{-4\sigma})}{32}
 -\sum_{p=3,5}\frac{\log(1-p^{-2\sigma})}{4}\\
 &-\sum_{p=7,11,13}\frac{\log(1-p^{-4\sigma})}{8}
 -\frac{\log(1-17^{-4\sigma})}{4}.
 \label{an:bad-restored}
\end{align}
This restores the deleted factors using the local indices of $K$.
The indices of $N$ at $S'$ may be different.

First consider $p\leq10^4$, $p\notin S'$, with nonzero genus vector
$v_p$. Its square is detected in the full retained degree-two quotient
if some form in the twelve-dimensional central space has $q(v_p)=1$.
If all seven completed forms instead have $q(v_p)=0$, then $f_N(p)=2$
whereas $4\mid f_K(p)$. Define $\mathcal P_4$ by precisely these
conditions. It contains $45$ primes. Equation~\eqref{an:local-phi}
gives the saving
\begin{equation}\label{an:order-four-saving}
 S_4=\frac14\sum_{p\in\mathcal P_4}
                     \log\frac{1+p^{-2\sigma}}{1-p^{-2\sigma}}
 >0.00003724741189544274807446.
\end{equation}
Primes for which the completed field already has order four are
excluded from this sum.

The second correction uses primes that split in $E$ and in $N$,
but not in the full twelve-dimensional retained central field.
Here $f_N(p)=1$ and $f_K(p)$ is even, so the saving is
\begin{equation}\label{an:central-saving}
 S_{\rm cen}=\sum_{\substack{p\leq10^{12}\\
                  p\text{ split in }N,\ p\text{ nonsplit in the full field}}}
                     \operatorname{atanh}(p^{-\sigma}).
\end{equation}
The support has zero genus vector, so it is disjoint from
$\mathcal P_4$. Both supports exclude $S'$. No saving from an
uncomputed tail is claimed.

To compute~\eqref{an:central-saving}, enumerate the genus-split primes.
Their residues are the $180$ classes modulo $120120$ satisfying
$p\equiv1\pmod8$ and being nonzero squares modulo each of
$3,5,7,11,13$. Construct them by the Chinese remainder theorem.
For each class, sieve the numbers $120120k+r$ through $X=10^{12}$
by all primes up to $\sqrt X$, beginning at the square of each
sieving prime. Every surviving number above $13$ is prime, since
a composite would have a prime divisor at most its square root.
The classes are disjoint and none contains a prime dividing the modulus.

For a surviving prime, evaluate the $17$ norm radicands of
Section~\ref{ar:catalog} using exact modular square roots. The symbol of
$x+y\sqrt A$ is independent of the choice of root: its product with
the conjugate is $Bz^2$, a nonzero square at a genus-split prime.
All catalog $B$ are supported on $2,3,5,7,11,13$ and all catalog
$z$ are $1$ or $2$. Thus no prime in this census divides $Bz$,
and no zero residue or unhandled norm prime occurs in these raw
predicates.
The seven completed and twelve full-field product predicates in
Proposition~\ref{ar:central-fields} are then exact splitting tests.
Order the completed tests first and count only the first nonzero
predicate. A hit in the first seven tests is already accounted for by
$N$. A later hit belongs to~\eqref{an:central-saving}. A prime on
which every predicate vanishes is not counted. This proves that the
computed complement contains each eligible prime exactly once.

For outward summation, use disjoint bins $l<p\leq h$, beginning at
$l=13$ and with next endpoint
$h=\min(X,\max(l+1,\lceil1001l/1000\rceil))$.
For complement primes in a bin record their exact number $n$ and
$u=\sum\lfloor10^{30}/p\rfloor$. With $\varepsilon=\sigma-1$,
their contribution is enclosed between
\begin{equation}\label{an:census-bin}
 h^{-\varepsilon}\frac{u}{10^{30}}
 \quad\hbox{and}\quad
 l^{-\varepsilon}\frac{u+n}{10^{30}}
                         \frac1{1-l^{-2}}.
\end{equation}
Indeed $x\leq\operatorname{atanh}x\leq x/(1-x^2)$ for $0<x<1$,
and $\sigma>1$. The only rounded quantities in the producer are the
explicit floors, whose full error is included in $u+n$.
All counts and reciprocal numerators use exact integers.
For the fixed-width census implementation, counts are at most $X$, and
the total reciprocal numerator is less than
$10^{30}(1+\log X)<2^{105}$. Thus its $64$-bit counts and $128$-bit
accumulators cannot overflow at this cutoff. Modular products use the
wider reduction described above.

The full sieve finds $293772662$ genus-split primes. Of these,
$293705608$ are detected by the full central field, and $291485548$
are already detected by $N$. The complement contains
$2220060$ primes. Summing~\eqref{an:census-bin} outward gives
\begin{equation}\label{an:census-value}
 0.00004201663320373563276702<S_{\rm cen}
          <0.00004201663670363397482805.
\end{equation}
An independent prime-list and modular-arithmetic calculation through
$10^7$ checks every bin, first-hit count, and reciprocal numerator.
Its complement has $31$ primes. This finite comparison checks the
implementation but does not replace the full sieve or its completeness
argument.

Primewise comparison now proves, for every retained tower field,
\begin{equation}\label{an:full-euler}
 \frac{\log\zeta_K(\sigma)}{[K:\mathbb Q]}
 \leq Y(\sigma)+B_D(\sigma)-S_4-S_{\rm cen}.
\end{equation}

\begin{remark}[Fixed-field form]\label{an:fixed-field-form}
The same primewise comparison applies to the retained field $M$ of
Proposition~\ref{ar:central-fields} itself, of degree $2^{19}$, and
gives
\[
 \frac{\log\zeta_M(\sigma)}{2^{19}}
 \leq Y(\sigma)+B_D(\sigma)-S_4-S_{\rm cen}.
\]
Since $M$ contains $N$, equation~\eqref{an:local-phi} bounds its
normalized zeta logarithm prime by prime by that of $N$ outside $S'$.
At $S'$ we restore the factors in $B_D$ using the indices of $M$.
These are $(e,f)=(8,4)$ at $2$, $(2,2)$ at $3$ and $5$, $(2,4)$ at
$7,11,13$, and $f=4$ at $17,19,23,29,31$. To justify these indices,
Proposition~\ref{ar:central-fields} identifies $M$ with the field of
$G/D_3G$, which retains every local cut of Section~\ref{tw:tower}.
A direct local computation from the radicands also confirms them.
They give the terms in \eqref{an:bad-restored} and the local floors
used below.

The conditions defining $\mathcal P_4$ and the census complement are
the tests \eqref{ar:nonsplit-test} on the central quadratic forms and
the twelve full-field predicates. These are tests in $M$. Hence a prime
in $\mathcal P_4$ has $f_M(p)=4$, and a complement prime has $f_M(p)=2$.
Both savings $S_4$ and $S_{\rm cen}$ are thus realized in $M$.
Each term in the following absolutely convergent series decreases in
$e_M(p)$ and $f_M(p)$, so the same local floors give
\[
 -\frac{\zeta_M'}{\zeta_M}(2)\cdot\frac1{2^{19}}
 =\sum_p\frac{\log p}{e_M(p)\,(p^{2f_M(p)}-1)}\leq R
\]
with $R$ as in \eqref{an:R}.
Consequently the fixed-field quantity
\[
 H_M:=\frac{\log\zeta_M(\sigma)}{2^{19}}
 +(\sigma-1)\left(\frac{\ell-\gamma_E-\log(4\pi)}4
 -\frac{\zeta_M'}{\zeta_M}(2)\cdot\frac1{2^{19}}\right)
\]
satisfies, with $L$ as in \eqref{an:slope} and the five allowances
\eqref{an:Y-numerical}, \eqref{an:bad-value},
\eqref{an:order-four-saving}, \eqref{an:census-value} and
\eqref{an:debit},
\[
 H_M\leq Y(\sigma)+B_D(\sigma)-S_4-S_{\rm cen}+(\sigma-1)L
 <0.04216181893948094953138458.
\]
The accompanying Lean formalization assumes this inequality for $H_M$
with the relaxed threshold $0.042165819$. Under that hypothesis, it
proves the transfer from $M$ to the tower fields and
Proposition~\ref{an:prime-budget}. The numerical bound for $H_M$ itself
is not proved in the formalization.
\end{remark}

\subsection{The Numerical Prime-Budget Bound}

We specify the local floors used in Proposition~\ref{an:prime-budget}.
For $p\leq10^4$, take $e_2^0=8$, $e_p^0=2$ at
$p=3,5,7,11,13$, and $e_p^0=1$ elsewhere. Use the exact retained
residue degrees $4$ at $2$, $2$ at $3,5$, and $4$ at
$7,11,13,17,19,23,29,31$. At any other such prime use $f_p^0=4$
if its square is nonzero in the full central quotient, $f_p^0=2$
if its genus vector is nonzero, and $f_p^0=1$ otherwise.
We apply the square test only when the genus vector is nonzero. For a
genus-split prime with a nontrivial central image, we keep the weaker
floor one. All these are lower bounds supplied by the
actual finite quotients in Proposition~\ref{ar:central-fields}.
Above $10^4$ use both floors one.

Since $\log x/(x^2-1)$ decreases on $[2,\infty)$, the integer tail
is bounded by its integral. It is enough to take
\begin{equation}\label{an:R}
 R=\sum_{p\leq10^4}\frac{\log p}{e_p^0(p^{2f_p^0}-1)}
          +\frac{\log10^4+1}{10^4(1-10^{-8})}.
\end{equation}
All finite floors are reconstructed by exact genus and square tests.
Outward evaluation gives
\begin{align}
 B_D(\sigma)&<0.041727023150898786,\label{an:bad-value}\\
 L:=\frac{\ell-\gamma_E-\log(4\pi)}4+R
 &<0.824453118933538946712643997041,\label{an:slope}\\
 (\sigma-1)L&<0.00006870442657779491222606.
 \label{an:debit}
\end{align}
The slope is independent of the abscissa, and the debit uses exactly
$\sigma-1=1/12000$.

Combining Proposition~\ref{an:prime-budget} with
\eqref{an:full-euler}, and then using the five conservative rational
bounds \eqref{an:Y-numerical}, \eqref{an:bad-value},
\eqref{an:order-four-saving}, \eqref{an:census-value}, and
\eqref{an:debit}, gives
\begin{align*}
 \limsup_{d\to\infty}\frac{\log L_F(1)}d
 &\leq Y(\sigma)+B_D(\sigma)-S_4-S_{\rm cen}+(\sigma-1)L\\
 &<0.04216181893948094953138458\\
 &<0.042161819.
\end{align*}
The middle number is the exact sum of the five rational bounds, with
the saving terms subtracted using their lower endpoints. The eventual
uniformity in Proposition~\ref{an:prime-budget} now proves
Theorem~\ref{an:ceiling}.

\section{Norm-One Mass and the Geometric Transfer}
\label{geo:section}

Let $K/F$ be quadratic with involution $\iota$, let $K$ be totally
imaginary, and let $F$ have signature $(b,c)$, where $b>0$. We keep these
assumptions in Sections~\ref{fw:section} and~\ref{pf:section}.
Write
\[
 d=[F:\mathbb Q]=b+2c,\qquad \theta=c/d,
 \qquad |\Delta_K|=\lambda^{2d},\quad
 |\Delta_F|=\lambda^d,\quad \ell=\log\lambda.
\]
Assume that $K/F$ is unramified at all finite places. The discriminants
are field discriminants, and archimedean integrals use Lebesgue measure
on unweighted complex coordinates. We count ordered edges and divide
by two to obtain the stated bounds for unordered edges.

We use Dirichlet's unit theorem, Kronecker's theorem, trace duality,
Hilbert's theorem~90, and the analytic class-number formula
\cite{MilneANT,MilneCFT}.

\subsection{The Unit Mass}

Put $E_K=\mathcal O_K^\times$, $E_F=\mathcal O_F^\times$,
$U=\ker(N:E_K\to E_F)$, and $I_U=[E_F:N(E_K)]$.
At each complex place of $F$, choose its two extensions to $K$ and set
\[
 l_j(x)=\frac12\bigl(\log|\sigma_j^+(x)|-
                         \log|\sigma_j^-(x)|\bigr).
\]
On $U$, this is $\log|\sigma_j^+(x)|$. Its image is a full lattice in
$\mathbb R^c$ of covolume $R_U$, where $R_U=1$ if $c=0$.
Let $w_K$ be the number of roots of unity in $K$,
$h_{\rm rel}=h_K/h_F$, $\kappa=|\ker(\mathrm{Cl}(F)\to\mathrm{Cl}(K))|$,
and $L_F(s)=\zeta_K(s)/\zeta_F(s)$. The following identities give the
mass of the norm-one units needed in the lattice construction.

\begin{proposition}\label{geo:mass}
We have that
\[
 \kappa=I_U/2^{b-1},\qquad
 R_K/R_F=2^{c-1}I_U R_U,
\]
and consequently
\begin{equation}\label{geo:mass-density}
 \frac{w_K}{h_{\rm rel}\kappa R_U}
 =\frac{2^{d-1}\pi^{b+c}}{\lambda^{d/2}L_F(1)}.
\end{equation}
Here $R_K,R_F$ are ordinary regulators, using twice the logarithm of
absolute value at a complex place.
\end{proposition}

\begin{proof}
Dirichlet's theorem gives invariant and anti-invariant rational unit
ranks $b+c-1$ and $c$. The Herbrand quotient is multiplicative in exact
sequences, equals one for finite groups, and is $2$ or $1/2$ on a trivial
or sign copy of $\mathbb Z$. Applying this to the finite-index sum of
the two eigenspace lattices gives the quotient $2^{b-1}$, without
requiring the integral unit lattice to split.
Its degree-zero Tate group has size $I_U$, so
$|H^1(\langle\iota\rangle,E_K)|=I_U/2^{b-1}$.

Capitulation injects into this $H^1$, since a relation
$\mathfrak a\mathcal O_K=(x)$ gives the cocycle $x/\iota x$ and a
coboundary gives an invariant generator. Conversely, Hilbert 90
writes every norm-one unit as $x/\iota x$. The fractional ideal $(x)$ is
invariant. It descends from $F$, since exponents agree on split prime
pairs and every finite ramification index is one. This proves
surjectivity and the formula for $\kappa$.

The norm map on rational unit spaces is onto the invariant space, since
the norm of an $F$-unit is its square. Its kernel has rank $c$.
The logarithmic kernel in $U$ consists of roots of unity by Kronecker's
theorem. Every root of unity of $K$ has norm one, since its norm is a
torsion unit of $F$ and is positive at a real place.
Thus the logarithmic kernel has exactly $w_K$ elements.

For the regulator identity, delete a coordinate over a real place of
$F$ in both weighted logarithmic unit spaces. The norm adds the two
coordinates above a complex place of $F$. Replacing a pair $(x,y)$ by
$(x,x+y)$ has determinant one. On the kernel its first coordinate is
$2l_j$, giving covolume $2^cR_U$. The logarithmic norm image has index
$I_U/2$ in the logarithmic $F$-unit lattice. The torsion index is two,
since $-1$ is not a norm at a complexifying real place. Multiplication
of the kernel and image covolumes gives
$R_K/R_F=2^cR_U(I_U/2)$.

The analytic class-number formula
\cite[Chapter~V, Theorem~2.4]{MilneCFT}, with $w_F=2$, gives
\[
 h_{\rm rel}\frac{R_K}{R_F}
 =\frac{w_K\lambda^{d/2}L_F(1)}{2^{c+1}\pi^{b+c}}.
\]
Substitution proves \eqref{geo:mass-density}.
\end{proof}

Choose finitely many primes $\mathfrak p$ of $F$ split as
$\mathfrak P\iota\mathfrak P$ in $K$, and integers $k_{\mathfrak p}\ge0$.
Set $P_0=\prod_{\mathfrak p}(k_{\mathfrak p}+1)$. The one-sided ideals
$B_{\boldsymbol j}=\prod\mathfrak P^{j_{\mathfrak p}}$ have a fiber of
size at least $P_0/(\kappa h_{\rm rel})$ in
$\mathrm{Cl}(K)/\mathrm{im}\,\mathrm{Cl}(F)$. If
$B_{\boldsymbol j}/B_{\boldsymbol j_0}=(x)\mathfrak a\mathcal O_K$, then
\[
 J_{\boldsymbol j}/J_{\boldsymbol j_0}=(x/\iota x),\qquad
 J_{\boldsymbol j}=\prod\mathfrak P^{j_{\mathfrak p}}
                (\iota\mathfrak P)^{k_{\mathfrak p}-j_{\mathfrak p}}.
\]
Thus $I=J_{\boldsymbol j_0}^{-1}$ contains this many disjoint cosets
$\beta_jU$ of norm-one elements. Their principal ideals distinguish the
cosets, and
\[
 N_K(I)=\prod_{\mathfrak p}(N_F\mathfrak p)^{-k_{\mathfrak p}}.
\]
For uniform rational-prime data define
\begin{equation}\label{geo:HJ}
 H=\sum_p\frac{k_p\log p}{e_p},\qquad
 J=\sum_p\frac{\log(k_p+1)}{e_pf_p}.
\end{equation}
Here $e_p,f_p$ are the absolute indices at the selected primes, and all
$F$-primes above them split in $K/F$. Then $N_K(I)=e^{-dH}$ and
$P_0=e^{dJ}$. Prime-ideal mixtures give the corresponding sums with a
bounded rounding error, provided the finite list of exponents is fixed.

\subsection{Unit Averaging and Translation}

For $h\in\mathbb R^c$, rescale each complex pair by $(e^{-h_j},e^{h_j})$,
leaving the $b$ compact coordinates unchanged. Denote the resulting ideal
lattice by $\Lambda_h$. Its determinant, independently of $h$, is
\begin{equation}\label{geo:determinant}
 D=2^{-d}\lambda^d e^{-dH}.
\end{equation}
For any nonnegative integrable $\Phi$ on $\mathbb R^c$, averaging over a
fundamental parallelepiped $P_U$ of $l(U)$ gives exactly
\begin{equation}\label{geo:unfold}
 \frac1{R_U}\int_{P_U}
 \sum_{\beta\in\bigcup_j\beta_jU}\Phi(l(\beta)-h)\,dh
 =\frac{q w_K}{R_U}\int_{\mathbb R^c}\Phi(u)\,du,
 \qquad q\ge\frac{P_0}{\kappa h_{\rm rel}}.
\end{equation}
This follows from tiling and Tonelli's theorem for the full unit lattice.
We use logarithmic Haar measure $du$ and angular measures of total mass
one.

Let $\Omega$ be bounded and invariant under rotations in each complex
coordinate. Let $V$ be its volume and let $\Phi(u)$ be its overlap with
a displacement having compact moduli one and pair moduli $(e^{u_j},e^{-u_j})$.
This function is nonnegative, bounded, and compactly supported in $u$.
Parameterize additive translates by $B_hr$, $r\in[0,1)^{2d}$, where
$B_h$ is a continuous basis matrix of $\Lambda_h$. The point count $N(h,r)$
and inherited ordered-edge count $E(h,r)$ are measurable countable sums
of indicators. Averaging $r$ by lattice tiling gives
\[
 \mathbb E_r N(h,r)=V/D,\qquad
 \mathbb E_r E(h,r)=D^{-1}\sum_\beta\Phi(l(\beta)-h).
\]
If a uniform point bound $N\le N_{\max}$ holds for every $h,r$, all these
quantities are integrable. Averaging also in $h$, and weighting parameter
space by $N$, selects one pair $(h,r)$ with $N>0$ and
\begin{equation}\label{geo:joint-average}
 \frac{E(h,r)}{N(h,r)}\ge
 \frac{q w_K}{R_U}\frac{\int\Phi}{V}.
\end{equation}
The vertex bound and the edge bound hold for the same rescaling and
translate.

\subsection{Uniform Lattice Counting}

The real Euclidean dual of the unweighted complex ideal lattice is the
coordinatewise complex conjugate of the image of
$2\mathfrak D_K^{-1}I^{-1}$. Indeed its real pairing is the trace pairing
divided by two. Reciprocal deformations invert on the dual. Consequently
every nonzero dual vector satisfies
\begin{equation}\label{geo:dual-product}
 \prod_{i=1}^d|\xi_i|\ge\mu^d,
 \qquad \mu=2e^{H/2-\ell},\qquad
 \sum_i|\xi_i|\ge d\mu.
\end{equation}
The product assertion follows by taking the square root of the absolute
norm and using $|N_K(x)|\ge1/(|\Delta_K|N_K(I))$ for
$0\ne x\in\mathfrak D_K^{-1}I^{-1}$.
Here conjugation acts coordinatewise and need not agree with $\iota$.

Suppose a smooth positive weight $F$ on $\mathbb C^d$, formed as a product
over the chosen one- and two-coordinate blocks, has normalized
Fourier bound, for the convention $e^{-2\pi i x\cdot\xi}$,
\[
 |\widehat F(\xi)|/\int F\le M^d
                  \exp\left(-\sigma\sum_i|\xi_i|\right),
 \qquad M\ge1,\quad \sigma>0.
\]
Put $R=d\mu$ and use the norm $\|\xi\|_*=\sum_i|\xi_i|$.
Every two distinct dual points are separated by at least $R$, so open
balls of radius $R/2$ centered at them are disjoint. Comparing volumes
inside the ball of radius $(j+3/2)R$ bounds the number of points in
$jR\le\|\xi\|_*<(j+1)R$ by
$(2j+3)^{2d}\le5^{2dj}$. Thus the sufficient condition
\begin{equation}\label{geo:poisson-condition}
 \sigma\mu\ge\log M+2\log5+2
\end{equation}
makes the nonzero normalized Fourier sum at most
$\sum_{j\ge1}e^{-2dj}<1$.
For Gaussian and polynomial Student profiles, we obtain
\begin{equation}\label{geo:pointwise-poisson}
 \sum_{x\in\Lambda_h+r}F(x)\le2\int F/D,
\end{equation}
uniformly in $h,r$. To justify Poisson summation, multiply $F$ by
$e^{-\varepsilon|x|^2}$. The result is Schwartz. Convolution with the
normalized positive Fourier Gaussian multiplies the Fourier envelope by
its exponential norm moment, which tends to one in each fixed dimension.
The strict tail bound persists for sufficiently small $\varepsilon$.
Schwartz Poisson summation and monotone convergence on the primal side
prove \eqref{geo:pointwise-poisson}. This limit is taken before $d$ grows.

For $F_C(z)=e^{-a_C|z|^2}$ with $0<a_C\le1$, one may use
$2e^{-|\xi|}$ as a normalized Fourier envelope. For
$F_C(z)=(1+a|z|^2)^{-\nu}$, $\nu>1$, Gaussian integration gives
\[
 \frac{\widehat F_C(\xi)}{\int F_C}
 =\frac1{\Gamma(\nu-1)}\int_0^\infty
 t^{\nu-2}e^{-t-\pi^2|\xi|^2/(at)}\,dt
 \le2^{\nu-1}e^{-\sqrt2\pi|\xi|/\sqrt a}.
\]
The inequality follows from $t+L/t\ge t/2+\sqrt{2L}$.

\subsection{Compact and Complex-Pair Profiles}

Fix $0<\delta<1$, $p=2/(1+\delta)$, a positive radial $f_C$ on $\mathbb C$,
and a positive separately radial $g_D$ on $\mathbb C^2$. Here $p$ denotes
the Lebesgue exponent. Define
\[
 A_C=\int f_C^p,\quad I_C=\int f_C(z)f_C(z+1)\,dz,
 \qquad A_D=\int g_D^p,
\]
\[
 I_D=\int_{\mathbb R}\int_{\mathbb C^2}
 g_D(x,y)g_D(x+e^u,y+e^{-u})\,dx\,dy\,du,
\]
\[
 J_C=\log I_C-(1+\delta)\log A_C,
 \quad J_D=\log I_D-(1+\delta)\log A_D.
\]
Assume these integrals are finite and positive, the two endpoint log-energy
variables have finite second moments, the total energy is bounded below
and coercive, and \eqref{geo:pointwise-poisson} holds for the tensor
weight $F=\prod f_C^p\prod g_D^p$ uniformly in reciprocal deformations.
Section~\ref{pf:section} verifies these hypotheses for our profiles. They
give the following geometric criterion.

\begin{proposition}\label{geo:transfer}
Suppose the preceding arithmetic data exist with $d\to\infty$, and
$\log L_F(1)\le dC$. With $H,J,\ell$ fixed up to bounded total rounding
errors, set
\begin{equation}\label{geo:general-margin}
 \mathcal M(\theta)=J-\delta H-(1/2-\delta)\ell-C
 +(1-\delta)\log2+(1-\theta)\log\pi
 +(1-2\theta)J_C+\theta J_D.
\end{equation}
If this margin has a fixed positive lower bound over all occurring
signatures, there are arbitrarily large planar sets whose unordered
unit-distance count divided by $|U|^{1+\delta}$ tends to infinity.
\end{proposition}

\begin{proof}
Write $V_C=-\log f_C$, $V_D=-\log g_D$. The compact endpoint law has
density $f_C(z)f_C(z+1)/I_C$ with respect to $dz$. The pair endpoint law
has density $g_D(x,y)g_D(x+e^u,y+e^{-u})/I_D$ with respect to $dx\,dy\,du$.
These are probability laws. Reflection about the midpoint exchanges the
two endpoints without changing either law. Let the first-endpoint energy
means be $v_C,v_D$ and set
$v=(b v_C+c v_D)/d$. Take the common sublevel set
\[
 \Omega=\left\{\sum V_C+\sum V_D\le d(v+\eta)\right\},\qquad\eta>0.
\]
It is bounded. Its volume is at most
$A_C^bA_D^c e^{pd(v+\eta)}$, and the pointwise weighted bound gives
\[
 N_{\max}\le2e^{dB_\eta},\quad
 B_\eta=H+\log2-\ell+(1-2\theta)\log A_C+\theta\log A_D+pv+p\eta.
\]
Independence across blocks, endpoint exchange symmetry, Chebyshev's inequality and a
union bound put both endpoint energies within $\eta d$ of their means
with probability at least $1/2$ for large $d$, uniformly in signature.
Undoing the overlap density gives
$\int\Phi\ge\tfrac12 I_C^b I_D^c e^{2d(v-\eta)}$.
Equations \eqref{geo:mass-density} and \eqref{geo:joint-average}
give ordered average degree at least $\tfrac14e^{dA_\eta}$, where
\[
 \begin{aligned}
 A_\eta={}&J+\log2+(1-\theta)\log\pi-\ell/2-C
 +(1-2\theta)\log I_C+\theta\log I_D\\
 &-(1-2\theta)\log A_C-\theta\log A_D
 +(2-p)v-(2+p)\eta.
 \end{aligned}
\]
Since $p(1+\delta)=2$, we have
$A_\eta-\delta B_\eta=\mathcal M(\theta)-4\eta$.
The ordered ratio is at least
$2^{-2-\delta}\exp(d(\mathcal M-4\eta))$, apart from bounded rounding
factors. Choose $4\eta<\inf\mathcal M$.
A compact coordinate comes from a real place of $F$. Projection to that
embedding is injective on differences of affine ideal points and sends
all distinguished norm-one displacements to unit vectors. Division by two
gives the unordered conclusion. Since $u(U)\le |U|^2/2$ and
$0<\delta<1$, the ratio $u(U)/|U|^{1+\delta}$ can tend to infinity only
if $|U|\to\infty$. This gives a sequence of cardinalities, with no claim
at every sufficiently large cardinality.
\end{proof}

\section{Weighted Windows at Split Finite Places}
\label{fw:section}

We extend the geometric criterion by placing locally constant weights
at the selected finite primes. Keep the notation $K/F$, $d=b+2c$,
$\theta=c/d$, $\lambda$ and $\ell=\log\lambda$ of
Section~\ref{geo:section}, where $b>0$ and $K/F$ is unramified at all
finite places. Put $D_0=(\lambda/2)^d$, and keep the archimedean profiles
$f_C,g_D$ and functionals $J_C,J_D$ defined there.

\subsection{The Local Functional}

Let $L$ be a nonarchimedean local field with valuation ring $\mathcal O$,
uniformizer $\pi$, and residue cardinality $Q$. Additive Haar measure gives
$\mathcal O^2$ mass one. Multiplicative Haar measure $du$ on
$\mathcal O^\times$ has mass one, and $n\in\mathbb Z$ has counting measure.
Define
\[
 s(n,u)=(\pi^nu,\pi^{-n}u^{-1}),\qquad
 Tg(z)=\sum_{n\in\mathbb Z}\int_{\mathcal O^\times}g(z+s(n,u))\,du.
\]
For a nonnegative, locally constant, compactly supported $g$, put
\begin{equation}\label{fw:functional}
 A(g)=\int g^p,\qquad Z(g)=\int gTg,\qquad
 \mathcal F_\delta(g)=\frac{Z(g)}{A(g)^{1+\delta}},
 \qquad p=\frac2{1+\delta},\quad0<\delta<1.
\end{equation}
Assume $g\ne0$ and $Z(g)>0$. Only finitely many valuation labels
contribute, since both coordinates of $s(n,u)$ lie in a fixed compact
difference set. The unit average has mass one.

For any integer $a$, put $B_a=\pi^{-a}\mathcal O$. Ball intersection gives
\begin{equation}\label{fw:rectangle-gram}
 \int 1_{B_a\times B_b}\,T1_{B_c\times B_e}
 =Q^{\min(a,c)+\min(b,e)}
   [\max(a,c)+\max(b,e)+1]_+.
\end{equation}
Two translated balls intersect with the smaller volume exactly when
the displacement lies in the larger ball. The first displacement must
have valuation at least $-\max(a,c)$ and the second at least
$-\max(b,e)$. These conditions say
$-\max(a,c)\le n\le\max(b,e)$. Every permitted unit contributes the same
intersection volume. This also proves the formula for negative ball indices.

Put $S_0=\mathcal O$, $S_i=B_i\setminus B_{i-1}$ for $i\ge1$, and
$m_0=1$, $m_i=Q^i-Q^{i-1}$. For $k\ge0$, let $g$ have value $w_{ij}$
on $S_i\times\pi^{-k}S_j$ for $0\le i\le L_x$, $0\le j\le L_y$, and
be zero elsewhere. The weights are nonnegative. With
\[
 D_{ir}=\begin{cases}
 m_{\min(i,r)},&i\ne r,\\
 0,&i=r=0,\\
 i m_i-Q^{i-1},&i=r\ge1,
 \end{cases}
\]
define
\[
 K_{(i,j),(r,t)}=(k+1)m_im_j\mathbf1_{i=r,\,j=t}
       +m_jD_{ir}\mathbf1_{j=t}+m_iD_{jt}\mathbf1_{i=r}.
\]
Then
\begin{equation}\label{fw:shell-kernel}
 A(g)=Q^k\sum_{i,j}m_im_jw_{ij}^p,\qquad Z(g)=Q^k w^TKw.
\end{equation}
To verify the second formula, take the four shell differences of
\eqref{fw:rectangle-gram}. The constant $k+1$ gives the diagonal
term. The double difference of $\max(i,r)Q^{\min(i,r)}$ is $D_{ir}$,
whereas that of $Q^{\min(j,t)}$ is $m_j\mathbf1_{j=t}$. The remaining
term is symmetric. Since $k\ge0$, the positive-part sign can be removed
in all these differences. The formula holds in residue characteristic
two and for ramified $L/\mathbb Q_p$.

These profiles are separately invariant under multiplication by units and
have the additive period
\begin{equation}\label{fw:hard-period}
 C=\mathcal O\times\pi^{-k}\mathcal O,\qquad |C|=Q^k,
\end{equation}
irrespective of the number of shells. In particular $w_{00}=1$ and all
other weights zero recover $\mathcal F_\delta=(k+1)Q^{-\delta k}$.

For product weights $w_{ij}=a_i b_j$, define
\[
 P_a=\sum_i m_i a_i^p,\quad S_a=\sum_i m_i a_i^2,\quad
 R_a=\sum_{i,r}a_iD_{ir}a_r,
\]
and define $P_b,S_b,R_b$ similarly. Formula
\eqref{fw:shell-kernel} becomes
\begin{equation}\label{fw:product-shells}
 A(g)=Q^kP_aP_b,\qquad
 Z(g)=Q^k[(k+1)S_aS_b+R_aS_b+S_aR_b].
\end{equation}
Thus symmetric product weights need only one vector. The one-shell choice
$(a_0,a_1)=(b_0,b_1)=(1,x)$ has
$S=1+(Q-1)x^2$ and $R=2x+(Q-2)x^2$, so its log gain over the hard
rectangle is
\begin{equation}\label{fw:one-shell-gain}
 \log\frac{(k+1)S^2+2SR}{k+1}
 -2(1+\delta)\log(1+(Q-1)x^p).
\end{equation}
Its right derivative at zero is $4/(k+1)>0$. For the numerical
certificate we use the full finite functional.

\subsection{Finite Periods and Lattice Counting}

Let $S$ now denote a finite set of $F$-primes split in $K$. At each
$v\in S$ identify the two completions with $L_v=F_v$, compatibly with
$\iota$, and choose a common uniformizer. Put
\[
 E_S=\prod_{v\in S}(L_v\times L_v),\qquad
 \Gamma=\mathcal O_{K,S}\subset\mathbb C^d\times E_S.
\]
Here $\mathcal O_{K,S}$ consists of the elements integral outside the
selected prime pairs. For a rectangular compact open subgroup
$C_f=\prod_v(\pi_v^{a_v}\mathcal O_v\times\pi_v^{b_v}\mathcal O_v)$,
the elements of $\Gamma$ with finite component in $C_f$ form the fractional
ideal
\[
 L=\prod_v\mathfrak P_v^{a_v}(\iota\mathfrak P_v)^{b_v},\qquad
 N_K(L)=|C_f|^{-1},\qquad \operatorname{covol}(L)=D_0/|C_f|.
\]
Every $C_f$-coset in $E_S$ contains a finite component of a $\Gamma$-point.
To see this, enclose the coset and $C_f$ in larger local balls and form
the corresponding fractional ideal $J\supset L$. After clearing a common
denominator, the Chinese Remainder Theorem identifies $J/L$ with its
local quotients at the selected prime ideals. This requires no
principality hypothesis, and each finite coset has an archimedean fiber
equal to a translate of $L$.

If $P_L$ is an archimedean fundamental parallelepiped, then $P_L\times C_f$
is a fundamental domain for $\Gamma$, of Haar volume $D_0$. First subtract
a $\Gamma$-element to put the finite component in $C_f$, then subtract an
$L$-element to put the infinite component in $P_L$. The same argument
proves uniqueness and discreteness. Reciprocal archimedean deformation
$B_h$ gives $\Gamma_h=(B_h,\mathrm{id})\Gamma$ with the same covolume.

Write $H_C=d^{-1}\log|C_f|$. The codifferent argument in
\eqref{geo:dual-product} applied to $L$ gives the product minimum
$[2e^{H_C/2-\ell}]^d$ on every deformed dual lattice. Consequently the
archimedean Fourier condition sufficient for this section is
\begin{equation}\label{fw:fourier-condition}
 \sigma\,2e^{H_C/2-\ell}\ge\log M+2\log5+2,
\end{equation}
using the envelope and pointwise regularity of Section~\ref{geo:section}.
It gives a bound $2\int F_\infty/(D_0/|C_f|)$ for the sum of
$F_\infty=\prod f_C^p\prod g_D^p$ on every archimedean fiber.
Summing against any nonnegative, compactly supported, $C_f$-periodic
function $F_f$ yields
\begin{equation}\label{fw:weighted-count}
 \sum_{\gamma\in\Gamma_h+w}F_\infty(\gamma_\infty)F_f(\gamma_f)
 \le\frac2{D_0}\int F_\infty\int F_f.
\end{equation}
The values of $F_f$ on its finitely many nonzero cells sum to
$(\int F_f)/|C_f|$, so the estimate depends on the period volume.
Local reciprocal diagonal rescaling replaces $C_f$ by a rectangular
period of the same volume. Hence the bound remains uniform after such
rescalings, including those that do not preserve $\Gamma$.

\subsection{Class Selection and Unit Averaging}

Let $\Omega\subset\mathbb C^d\times E_S$ be bounded, have compact finite
support, and be invariant under all coordinate rotations and local
reciprocal unit multiplications. Denote its overlap with a norm-one step
of archimedean pair-log coordinates $u\in\mathbb R^c$ and finite valuation
labels $n\in\mathbb Z^S$ by $\Phi(u,n)$. Invariance removes all phase and
unit choices. Only finitely many labels occur, and the $u$-support is
bounded. Form the weights
\[
 W(n)=\int_{\mathbb R^c}\Phi(u,n)\,du.
\]
The quotient $\mathrm{Cl}(K)/\mathrm{im}\,\mathrm{Cl}(F)$ has order
$\kappa h_{\rm rel}$. Map $n$ to the class of
$A(n)=\prod_v\mathfrak P_v^{n_v}$. One class has total weight at least
$\sum_nW(n)/(\kappa h_{\rm rel})$. Suppose this sum is positive and
choose a positive-weight reference label $n_0$ in that class.
For every selected label, write
$A(n)/A(n_0)=(x_n)\mathfrak a_n\mathcal O_K$. Then
\[
 \beta_n=x_n/\iota(x_n),\quad N_{K/F}\beta_n=1,\quad
 (\beta_n)=\prod_v\mathfrak P_v^{n_v-n_{0,v}}
                  (\iota\mathfrak P_v)^{-n_v+n_{0,v}}.
\]
Then $\beta_n\in\mathcal O_{K,S}$ and the full set of norm-one generators
of this ideal is $\beta_nU$. Distinct labels give disjoint cosets.

Let $\alpha_0$ be the identity at infinity with finite coordinates
$(\pi_v^{n_{0,v}},\pi_v^{-n_{0,v}})$, and put $\Omega'=\alpha_0^{-1}\Omega$.
This preserves Haar and period volumes. The overlap with the displacement
$(B_h\beta,\beta_f)$, for $\beta\in\beta_nU$, is
$\Phi(l(\beta)-h,n)$. Since the overlap is independent of units and
phases, correlations between local units of global generators and their
archimedean data do not affect it. Tiling by the full unit lattice gives
\[
 \frac1{R_U}\int_{P_U}\sum_{n\text{ selected}}
 \sum_{\beta\in\beta_nU}\Phi(l(\beta)-h,n)\,dh
 =\frac{w_K}{R_U}\sum_{n\text{ selected}}W(n).
\]
For each fixed field and bounded $h$-domain, only finitely many unit-log
lattice points meet the bounded overlap support, and the logarithmic
kernel has $w_K$ elements. Tonelli also justifies the sums directly,
since their terms are nonnegative. We need no distribution hypothesis
on the local unit images.

Jointly average $h$ and additive translations over a fundamental domain
of $\Gamma_h$, then weight parameter space by its point count. The
translation averages are $|\Omega|/D_0$ for vertices and the sum of
overlaps divided by $D_0$ for ordered edges. Hence some pair of
transformations gives a nonempty point set satisfying
\begin{equation}\label{fw:joint-average}
 \frac EN\ge
 \frac{w_K}{\kappa h_{\rm rel}R_U}
 \frac{\sum_n\int\Phi(u,n)\,du}{|\Omega|}.
\end{equation}
The uniform bound \eqref{fw:weighted-count} holds for this same pair.
Choose the parameter domains using a continuously deformed basis and
Haar measure on the finite period. Countable indicator sums and Tonelli
justify the averages, including sublevel boundaries, since the
exponential majorant is pointwise and additive tiling holds for
measurable indicators. The class is chosen from the overlaps of the
final window, so concentration will be applied before class selection.

\subsection{The Finite-Place Transfer}

For every selected place choose a nonzero, nonnegative, locally constant,
compactly supported $g_v$ with a rectangular additive period $C_v$ and
the unit invariance above, and with $Z_v=Z(g_v)>0$. Suppose their number
is $O(d)$ and their profiles belong to a fixed finite list of local types.
The finite shell profiles above meet these hypotheses. Write
$A_v=A(g_v)$ and $\mathcal F_{\delta,v}=Z_v/A_v^{1+\delta}$. These local
functionals give the following extension of Proposition~\ref{geo:transfer}.

\begin{proposition}\label{fw:transfer}
Suppose the field degrees tend to infinity, $\log L_F(1)\le dC$, and
\eqref{fw:fourier-condition} holds with $C_f=\prod_v C_v$.
Assume the archimedean energy and moment hypotheses of
Proposition~\ref{geo:transfer}. If
\begin{equation}\label{fw:margin}
 \begin{split}
 \mathcal M_{\rm fin}(\theta)={}&
 \frac1d\sum_{v\in S}\log\mathcal F_{\delta,v}
 -(1/2-\delta)\ell-C+(1-\delta)\log2+(1-\theta)\log\pi\\
 &+(1-2\theta)J_C+\theta J_D
 \end{split}
\end{equation}
has a fixed positive lower bound, there are planar sets of cardinality
tending to infinity whose unordered unit-distance count divided by
$n^{1+\delta}$ tends to infinity.
\end{proposition}

\begin{proof}
The total profile $f=\prod f_C\prod g_D\prod_v g_v$ has
\[
 A=\int f^p=A_C^bA_D^c\prod_v A_v,\qquad
 Z=I_C^bI_D^c\prod_v Z_v.
\]
Use the independent normalized overlap probability laws of all these
blocks. The two endpoint energy distributions agree by reflection and
the stated symmetries. Let $d\bar v$ be the sum of first-endpoint means
for energy $-\log f$. On each finite factor, the positive support is a
finite union of additive cells with a positive minimum weight, so its
endpoint energy is bounded. The total variance is $O(d)$.

On the positive support put
$\Omega=\{-\log f\le d(\bar v+\varepsilon)\}$, for $\varepsilon>0$.
Coercivity makes it bounded. Its finite support is compact and it has
all the required periods and invariances. Exponential majorization and
\eqref{fw:weighted-count}, also after the reference-label rescaling,
give
\[
 |\Omega|\le A e^{pd(\bar v+\varepsilon)},\qquad
 N\le 2A e^{pd(\bar v+\varepsilon)}/D_0.
\]
Chebyshev's inequality and a union bound put both endpoint energies in
$[d(\bar v-\varepsilon),d(\bar v+\varepsilon)]$ with probability at least
$1/2$ for large $d$. Undoing the overlap density on this event gives
\[
 \sum_n\int\Phi(u,n)\,du\ge\tfrac12 Z e^{2d(\bar v-\varepsilon)}.
\]
Since $p(1+\delta)=2$, \eqref{fw:joint-average} implies
\[
 \frac E{N^{1+\delta}}\ge
 \frac{w_K D_0^\delta}{2^{1+\delta}\kappa h_{\rm rel}R_U}
 \frac{Z}{A^{1+\delta}}e^{-4\varepsilon d}.
\]
Substituting \eqref{geo:mass-density}, the ordered ratio is at least
$2^{-2-\delta}e^{d(\mathcal M_{\rm fin}-4\varepsilon)}$.
Choose $4\varepsilon<\inf\mathcal M_{\rm fin}$ and divide by two for
unordered edges. Projection to a compact coordinate is injective on the
field points and sends the norm-one displacements to unit vectors.
The upper bound $N^2/2$ for unordered edges forces $N\to\infty$.
\end{proof}

If a rational prime has actual indices $e_p,f_p$ and all its $F$-places
split in $K/F$, there are $d/(e_pf_p)$ such places. A common local profile
contributes $\log\mathcal F_{\delta,p}/(e_pf_p)$ to
\eqref{fw:margin}. Hard rectangles give
$\log\mathcal F_{\delta,p}=\log(k_p+1)-\delta k_pf_p\log p$, so their
sum is exactly $J-\delta H$. Adding finite shells with period
\eqref{fw:hard-period} preserves
$H_C=\sum_p k_p\log p/e_p=H$. Hence the finite-place functional replaces
the ideal-choice term and leaves the Fourier condition unchanged.

\section{Archimedean Profiles}
\label{pf:section}

For a pair of complex places, set
\[
 t=(1+a|z|^2)^{-1},\qquad u=(1+a|w|^2)^{-1},\qquad
 g_D(z,w)=t^su^sP(t,u),
\]
where
\[
 s=1.14600585,\qquad a=0.0000128633241758,
\]
and
\[
 P(t,u)=\sum_{i,j=0}^3 B_{ij}b_i^3(t)b_j^3(u),\qquad
 b_i^m(t)=\binom mi t^i(1-t)^{m-i},
\]
with
\begin{equation}\label{pf:bernstein-witness}
 B=\begin{pmatrix}
 1&2.4129438977&2.7721799210&6.6857798061\\
 2.4129438977&3.8532983888&4.6371750119&8.1307168183\\
 2.7721799210&4.6371750119&5.4628959084&9.5512647588\\
 6.6857798061&8.1307168183&9.5512647588&13.8659430610
 \end{pmatrix}.
\end{equation}
All finite decimals are exact rationals. The Bernstein basis is
nonnegative and sums to one, so $P$ is bounded above and below by
positive constants on the real square. We prove the hypotheses of
Proposition~\ref{fw:transfer} for these profiles and derive finite bounds
for their functionals. Section~\ref{sec:certificate} evaluates the
bounds at the fixed rational parameters.

\subsection{The Gaussian and the Pair Convolution}
The compact Gaussian $f_C^p=e^{-a_C|z|^2}$ gives
\begin{equation}\label{pf:gaussian-compact}
 J_C=\log(p/2)+\delta\log(a_C/\pi)-a_C/(2p),\qquad
 a_C=2p\delta
\end{equation}
where the displayed scale maximizes the Gaussian functional. For our
choice of $\delta$, we have $a_C<1$. The slope of the general margin as
a function of the signature
is $-\log\pi-2J_C+J_D$, and depends on the chosen profile.

For the pair calculation, let $\alpha,\alpha'>0$ with
$A=\alpha+\alpha'-1>0$. The Laplace identity
\[
 (1+|z|^2)^{-\alpha}
 =\frac1{\Gamma(\alpha)}\int_0^\infty
 r^{\alpha-1}e^{-r-r|z|^2}\,dr
\]
and Gaussian integration give
\begin{equation}\label{pf:laplace-convolution}
 \int_{\mathbb C}(1+|z|^2)^{-\alpha}
          (1+|z+h|^2)^{-\alpha'}\,dz
 =\frac\pi A\,
 \mathbb E(1+T(1-T)|h|^2)^{-A},
 \qquad T\sim\operatorname{Beta}(\alpha,\alpha').
\end{equation}
After integrating $z$, change the two positive Laplace variables to their
sum $r$ and ratio $T$. Completing the square gives the term
$rT(1-T)|h|^2$, and integration in $r$ gives the Beta expectation.
Tonelli justifies the interchanges, since the integrands are nonnegative.

For $A,B>0$ define
\[
 Q_{AB}(b)=\int_{\mathbb R}
 (1+be^{2v})^{-A}(1+be^{-2v})^{-B}\,dv,\qquad b>0.
\]
The two coordinates in the pair convolution reduce to this integral with
\[
 b=a\sqrt{T(1-T)T'(1-T')}\le a/4.
\]
This follows by translating $v$ until the coefficients of $e^{2v}$ and
$e^{-2v}$ agree, which has Jacobian one. The exponential representation is
\[
 \int_{\mathbb R}e^{-a(Re^{2v}+R'e^{-2v})}\,dv
 =K_0(2a\sqrt{RR'}),\qquad
 K_0(x)=\int_0^\infty e^{-x\cosh t}\,dt,
\]
where the substitution $t=2v+\frac12\log(R/R')$ and evenness cancel
their factors of two.

Splitting the $Q_{AB}$ integral at zero yields
\begin{equation}\label{pf:elementary-hyperbola-error}
 \left|Q_{AB}(b)-\log(1/b)
       +\frac{\psi(A)+\psi(B)+2\gamma_E}{2}\right|
 \le (A+B)b/2.
\end{equation}
Deleting the smaller factor on one half leaves
$\frac12\int_b^\infty(1+x)^{-A}dx/x$. Its difference from its
logarithmic constant lies between zero and $Ab/2$, by
$0\le1-(1+x)^{-A}\le Ax$. On that half the factor removed costs at most
$Bb/2$, since $1-(1+be^{-2v})^{-B}\le Bbe^{-2v}$.
The other half interchanges $A$ and $B$. Thus the total positive error
from the two tails and the total loss from the factors removed are both
at most $(A+B)b/2$, proving the bound. The logarithmic constant follows
from \cite[Equation~(5.9.16)]{DLMF},
\[
 \int_0^1\frac{(1-x)^{A-1}-1}{x}\,dx=-\psi(A)-\gamma_E.
\]
We use $\psi=\Gamma'/\Gamma$ and Euler's constant $\gamma_E$. For later
formulas write
\begin{equation}\label{pf:student-constant}
 C_s=-2\psi(s)+\psi(2s)+(2s-1)^{-1}-\gamma_E.
\end{equation}

\subsection{Endpoint Moments and the Fourier Bound}

The compact Gaussian has finite positive mass and overlap, and every
polynomial moment under its endpoint law is finite. Since $P$ is bounded
above and below by positive constants, the pair mass is finite for
$sp>1$. Formula~\eqref{pf:laplace-convolution} and
\eqref{pf:elementary-hyperbola-error} show that the pair overlap is finite:
the Beta variables have positive parameters and finite logarithmic
moments, and $b\le a/4$. Both integrals are positive. The energy
$-\log g_D=s\log(1+a|z|^2)+s\log(1+a|w|^2)-\log P(t,u)$ is bounded
below and coercive, so its sublevel sets are compact.

Finite endpoint-energy second moments hold for every positive polynomial
Student profile $t^su^sP(t,u)$ when $P$ is bounded above and below by
positive constants on the square and $sp>1$. Choose
$0<\epsilon<\min(s,2s-1)$ and absorb each logarithmic square into a
factor $(1+a|z|^2)^\epsilon$. Formula~\eqref{pf:laplace-convolution}
then has positive Beta parameters and positive hyperbola exponents.
The bound in \eqref{pf:elementary-hyperbola-error} is a constant plus
$|\log b|+b(A+B)/2$, which is integrable against the Beta densities.
This proves the second-moment condition, including $s=1$.

For the Fourier bound, write
$P(t,u)=\sum B_{ij}b_i^m(t)b_j^m(u)$ with $B_{ij}>0$. For a tube width
$0<\rho<1$, put $\omega=(\rho+\rho^2)/(1-\rho^2)$ and
\[
 E_{rs}=\binom mr\binom ms\max|\Delta_1^r\Delta_2^sB|,
 \qquad q_0=\frac{\sum_{r+s>0}E_{rs}\omega^{r+s}}{\min B}.
\]
If $q_0<1$, the Bernstein derivative formula puts $P$ in the right
half-plane throughout the product complex tube. Contour shifting then
gives, with $\beta=sp$,
\[
 |\widehat {g_D^p}(\xi_z,\xi_w)|/A_D
 \le (1-\rho^2)^{-2\beta}(1+q_0)^p
 e^{-2\pi\rho(|\xi_z|+|\xi_w|)/\sqrt a}.
\]
Here the tube is taken after the scaling $x=\sqrt a\,z$ in each
complex coordinate, and its imaginary displacement has Euclidean norm at
most $\rho$. If $D(x,y)=1+\sum_j(x_j+iy_j)^2$, then
\[
 |D(x,y)|\ge (1-\rho^2)(1+|x|^2),\qquad
 \left|D(x,y)^{-1}-(1+|x|^2)^{-1}\right|\le\omega.
\]
The latter follows by bounding the numerator by $2|x|\rho+\rho^2$
and the denominator by $(1-\rho^2)(1+|x|^2)^2$.
The Bernstein derivative identity bounds the coefficient of the
$(r,s)$ Taylor term by $E_{rs}$, uniformly at real $(t,u)\in[0,1]^2$.
Hence the complex value of $P$ differs from its positive real value by
at most $q_0P(t,u)$. We can define $P^p$ holomorphically in the tube and have
$|P_{\rm shifted}|^p\le(1+q_0)^pP(t,u)^p$. The displayed bound on $D$
supplies the other factor $(1-\rho^2)^{-2sp}$. Shift each real
Fourier contour by a vector of length $\rho/\sqrt a$ opposite its
frequency. The shifted majorant is integrable since $sp>1$, so
truncation and dominated convergence justify the shift. This proves the
Fourier envelope. The tube condition $q_0<1$, in addition to positivity
on the real square, gives explicit $M,\sigma$ in
\eqref{geo:poisson-condition}.

With the compact Gaussian \eqref{pf:gaussian-compact}, a sufficient
global choice in \eqref{geo:poisson-condition} is
\begin{equation}\label{pf:polynomial-fourier-contract}
 M=\max\{2,\sqrt{K_D}\},\quad
 \sigma=\min\{1,2\pi\rho/\sqrt a\},\quad
 K_D=(1-\rho^2)^{-2sp}(1+q_0)^p.
\end{equation}
There are $b$ one-coordinate blocks and $c$ two-coordinate blocks, so their
prefactors are bounded by $M^{b+2c}=M^d$. The maximum and minimum give
a bound uniform in the signature.

\subsection{Finite Bounds for Polynomial Profiles}

Let $s\ge1$, $\beta=sp>1$, $\eta_0=2s-1$, and
\[
 g_D(z,w)=t^s u^s P(t,u),\qquad
 P(t,u)=\sum_{i,j}d_{ij}t^iu^j>0\quad(0\le t,u\le1).
\]
Put
\[
 A_{ik}=\eta_0+i+k,\quad B_{jl}=\eta_0+j+l,\quad
 w_{ijkl}=\frac{d_{ij}d_{kl}}{A_{ik}B_{jl}},\quad b_0=\sum w_{ijkl}>0,
\]
\[
 c(x,y)=\frac1{x+y-1}
 +\frac{\psi(x+y-1)-\psi(x)-\psi(y)-\gamma_E}{2},
\]
\[
 \kappa_P=\frac{\sum w_{ijkl}[c(s+i,s+k)+c(s+j,s+l)]}{b_0}.
\]
Here $\pi^2b_0=\int_{\mathbb C^2}g_D^2$ at scale $a=1$.
Let $z_0=a^2/16$ and assume, for every cross with nonzero coefficient,
\begin{equation}\label{pf:cross-threshold}
 A_{ik}B_{jl}z_0<1,\qquad
 2\log(4/a)\ge H_{\lceil A_{ik}\rceil-1}+H_{\lceil B_{jl}\rceil-1},
 \qquad \log(4/a)>1/2,
\end{equation}
where $H_0=0$. Define
\[
 E_- =\log(4/a)\sum_{w_{ijkl}<0}(-w_{ijkl})
                \frac{A_{ik}B_{jl}z_0}{1-A_{ik}B_{jl}z_0}.
\]
Then
\begin{equation}\label{pf:polynomial-overlap}
 I_D\ge\pi^2a^{-2}b_0
       [\log(1/a)+\kappa_P-E_-/b_0].
\end{equation}
The bracket must be positive before taking its logarithm.

To prove this bound, recall that for $A,B\ge1$,
\[
 Q_{AB}(b)=\int_{\mathbb R}(1+be^{2u})^{-A}(1+be^{-2u})^{-B}\,du.
\]
For $0<b<1$, its convergent logarithmic series is
\[
 Q_{AB}(b)=\sum_{n\ge0}\frac{(A)_n(B)_n}{(n!)^2}b^{2n}
 \left[\log(1/b)+\psi(n+1)-\frac{\psi(A+n)+\psi(B+n)}2\right].
\]
To verify the series, substitute $x=be^{2u}$ and then $t=x/(x+b^2)$.
This expresses $2Q_{AB}(\sqrt z)$ as
\[
 \int_0^1t^{B-1}(1-t)^{A-1}(1-(1-z)t)^{-A}\,dt.
\]
Differentiation and integration by parts give
$z(1-z)y''+[1-(A+B+1)z]y'-AB y=0$.
Substituting the displayed series verifies its coefficients, using
$a_{n+1}/a_n=(n+A)(n+B)/(n+1)^2$ and the digamma recurrence.
The coefficients have polynomial growth, so the series and its fixed
derivatives converge locally uniformly for $0<z<1$. Its leading
logarithm and constant agree with the elementary estimate~\eqref{pf:elementary-hyperbola-error}. To see uniqueness with these two leading terms, let
$\phi(z)=\sum a_nz^n$, so $\phi(0)=1$ and $\phi>0$ on $[0,1)$.
For the difference $v$ of two such solutions, reduction of order gives
$(v/\phi)'=C/[z(1-z)^{A+B}\phi^2]$. Thus
$v/\phi=C\log z+D+o(1)$, and the matching leading terms force $C=D=0$.

The digamma recurrence and monotonicity imply
\[
 -H_{\lceil A\rceil-1}-H_{\lceil B\rceil-1}
 \le2\psi(n+1)-\psi(A+n)-\psi(B+n)\le0.
\]
Under the threshold in \eqref{pf:cross-threshold}, every term with
$n\ge1$ is nonnegative. Since $(A)_n(B)_n/(n!)^2\le(AB)^n$, summation gives
\[
 0\le Q_{AB}(b)-\log(1/b)
           +\tfrac12[\psi(A)+\psi(B)+2\gamma_E]
 \le\frac{ABb^2\log(1/b)}{1-ABb^2}.
\]
The right side increases for $b\le a/4$ under the stated conditions.
Apply this inside the two Beta expectations from \eqref{pf:laplace-convolution}, where $b=a\sqrt{T(1-T)T'(1-T')}\le a/4$.
Their logarithmic expectations give $c(x,y)$. For positive cross
coefficients we discard the nonnegative remainder. For negative cross
coefficients we subtract its upper bound. This proves
\eqref{pf:polynomial-overlap}, allowing either sign of monomial
coefficient.

To retain cancellation between monomial coefficients, we collect exact
terms before bounding the signed remainder. For $n\ge0$, define
\[
 K_n(i,k)=\frac{(s+i)_n(s+k)_n}
 {n!(\eta_0+i+k+n)_{n+1}},\qquad
 S_x(j)=\sum_{h=0}^{j-1}\frac1{x+h},
\]
with $S_x(0)=0$, and put
\[
 \begin{split}
 r_n(i,k)={}&-\eta_0^{-1}+H_n/2
 -S_{\eta_0}(i+k+n)/2+S_{\eta_0}(i+k+2n+1)\\
 &-[S_s(i+n)+S_s(k+n)]/2,\qquad R_n(i,k)=K_n(i,k)r_n(i,k),
 \end{split}
\]
\[
 B_n=\sum d_{ij}d_{kl}K_n(i,k)K_n(j,l),\qquad
 V_n=\sum d_{ij}d_{kl}
 [R_n(i,k)K_n(j,l)+K_n(i,k)R_n(j,l)].
\]
For rational parameters these are rational numbers. Under
\eqref{pf:cross-threshold}, for every integer $m\ge0$,
\begin{equation}\label{pf:collected-overlap}
 \begin{split}
 \frac{I_D}{\pi^2a^{-2}}\ge{}&
 \sum_{n=0}^m a^{2n}\{B_n[\log(1/a)+C_s]+V_n\}-E_{-,m},\\
 E_{-,m}={}&\log(4/a)\sum_{w_{ijkl}<0}(-w_{ijkl})
 \frac{(A_{ik}B_{jl}z_0)^{m+1}}{1-A_{ik}B_{jl}z_0}.
 \end{split}
\end{equation}
Thus $m=1$ retains the exact $a^2$ term and leaves a signed error of order
$a^4\log(1/a)$ for a fixed polynomial. The displayed right-hand side is a
lower bound for the overlap and gives no upper bound.

For a cross with exponents $\alpha=s+i$, $\alpha'=s+k$ and
$A=\alpha+\alpha'-1$, the coordinate factor at order $n$ is
\[
 \frac{(A)_n}{A n!}\,
 \mathbb E[(T(1-T))^n]
 =\frac{(\alpha)_n(\alpha')_n}
 {n!(A+n)_{n+1}}=K_n(i,k),\quad
 T\sim\mathrm{Beta}(\alpha,\alpha').
\]
The logarithmic moment uses the tilted distribution
$\mathrm{Beta}(\alpha+n,\alpha'+n)$, giving the coordinate constant
\[
 \tfrac12\psi(n+1)-\tfrac12\psi(A+n)
 -\tfrac12[\psi(\alpha+n)+\psi(\alpha'+n)]
 +\psi(A+2n+1)=C_s/2+r_n(i,k).
\]
The digamma recurrence proves the rational expression for $r_n$.
The remainder of each $Q_{AB}$ after index $m$ is nonnegative and bounded
by $\log(1/b)(ABb^2)^{m+1}/(1-ABb^2)$. Monotonicity for $b\le a/4$
under \eqref{pf:cross-threshold} proves \eqref{pf:collected-overlap}.
We also have $B_n>0$, since $K_n$ has the Gram representation
\[
 K_n(i,k)=\frac{(s+i)_n(s+k)_n}{(n!)^2}
 \int_0^1 t^{\eta_0+n-1}(1-t)^n t^i t^k\,dt.
\]
It is positive definite on any finite collection of distinct monomials,
and its tensor square makes $B_n$ the squared norm of a nonzero
polynomial. Individual signed crosses can still be negative.

For the mass, let $T,U$ be independent $\mathrm{Beta}(\beta-1,1)$ variables.
Radial substitution gives
\begin{equation}\label{pf:polynomial-mass}
 A_D=\frac{\pi^2a^{-2}}{(\beta-1)^2}\mathcal A,
 \qquad \mathcal A=\mathbb E P(T,U)^p.
\end{equation}
For a rational positive Bernstein representation, first clear a common
denominator. Since scalar multiplication of $P$ cancels from $J_D$, normalize
it to $P=1+V$ with $|V|\le r<1$ by dividing by half the sum of the smallest
and largest Bernstein coefficient. Then, for $N\ge2$,
\begin{equation}\label{pf:polynomial-mass-tail}
 \left|\mathcal A-\sum_{k=0}^N\binom pk\mathbb EV^k\right|
 \le\frac{|\binom p{N+1}|r^{N+1}}{1-r}.
\end{equation}
Indeed the binomial series is uniformly convergent on $[-r,r]$ and the
absolute binomial coefficients decrease after index two, since
$1<p<2$. If $\beta-1=A/C$ is positive rational, the monomial moments
$\mathbb ET^j=A/(A+Cj)$ make every coefficient in the truncated sum
exactly rational. We calculate them by integer polynomial convolution
and contraction. Since the mass has a negative logarithmic coefficient
in the functional, its upper endpoint gives a lower bound for the score.

For polynomials with a large ratio of maximum to minimum Bernstein
coefficient, the following local version has a smaller tail. Partition
$[0,1]^2$ into rational rectangles, put $q=\beta-1>0$, and on
$\mathcal R=[l,r]\times[b,t]$ use coordinates
$x=(T-l)/(r-l)$ and $y=(U-b)/(t-b)$. Restrict the polynomial exactly and
express it in the Bernstein basis in $x,y$. If its coefficient bounds are
$0<m_{\mathcal R}\le M_{\mathcal R}$, set
\[
 h_{\mathcal R}=(m_{\mathcal R}+M_{\mathcal R})/2,\quad
 \rho_{\mathcal R}=
 \frac{M_{\mathcal R}-m_{\mathcal R}}{M_{\mathcal R}+m_{\mathcal R}}<1,
 \quad P=h_{\mathcal R}(1+W_{\mathcal R}).
\]
The Bernstein convex-hull property gives
$|W_{\mathcal R}|\le\rho_{\mathcal R}$ throughout the rectangle.
Define the unnormalized one-dimensional moments
\begin{equation}\label{pf:subrectangle-moments}
 \begin{split}
 M_i(l,r)&=\int_l^r qv^{q-1}
       \left(\frac{v-l}{r-l}\right)^i\,dv\\
 &=\frac q{(r-l)^i}\sum_{j=0}^i\binom ij(-l)^{i-j}
             \frac{r^{q+j}-l^{q+j}}{q+j}.
 \end{split}
\end{equation}
The identity follows by expanding $(v-l)^i$. The convention $0^q=0$
is valid because $q>0$. If
$W_{\mathcal R}^k=\sum c^{(k)}_{ij}x^iy^j$, put
\[
 T_{\mathcal R,N}=h_{\mathcal R}^{p}
 \sum_{k=0}^N\binom pk\sum_{i,j}c^{(k)}_{ij}M_i(l,r)M_j(b,t).
\]
Then for $N\ge2$,
\begin{equation}\label{pf:subrectangle-mass}
 \left|\mathcal A-\sum_{\mathcal R}T_{\mathcal R,N}\right|
 \le\sum_{\mathcal R}h_{\mathcal R}^{p}
 \frac{|\binom p{N+1}|\rho_{\mathcal R}^{N+1}}
      {1-\rho_{\mathcal R}}M_0(l,r)M_0(b,t).
\end{equation}
This follows by integrating the uniform binomial-tail bound against the
nonnegative density on each rectangle. Polynomial restriction,
conversion, and powers have rational coefficients and may be calculated
by integer convolution after clearing denominators. The endpoint powers
in \eqref{pf:subrectangle-moments} and $h_{\mathcal R}^p$ are enclosed
using outward-rounded logarithms and exponentials, since they need not
be rational. Exact rational algebra followed by interval contraction
preserves the bound when the moment formula has cancellation. For symmetric $P$
and the same partition on both axes, off-diagonal rectangles may be
counted twice and diagonal rectangles once. Otherwise all rectangles
must be included separately.

Combining \eqref{pf:polynomial-overlap} and
\eqref{pf:polynomial-mass} gives
\begin{equation}\label{pf:polynomial-functional}
 \begin{split}
 J_D\ge{}&\log b_0+2(1+\delta)\log(\beta-1)-2\delta\log\pi
 -(1+\delta)\log\mathcal A+2\delta\log a\\
 &+\log[\log(1/a)+\kappa_P-E_-/b_0].
 \end{split}
\end{equation}
More generally, if $Z_*>0$ is the lower expression in
\eqref{pf:collected-overlap} and $\mathcal A^*$ is an upper bound from
\eqref{pf:subrectangle-mass}, we obtain the bound used in the certificate,
\begin{equation}\label{pf:collected-functional}
 J_D\ge\log Z_*+2(1+\delta)\log(\beta-1)-2\delta\log\pi
       +2\delta\log a-(1+\delta)\log\mathcal A^*.
\end{equation}
Interval evaluation encloses this lower-bound expression. Its upper
endpoint need not bound $J_D$.
For rational $s$ and rational polynomial data all digamma shifts in
$\kappa_P$ reduce to rational corrections and the single constant
\[
 C_s=\frac2{2s-1}-2(s-1)^2
 \sum_{k\ge1}\frac1{k(k+s-1)(k+2s-2)}.
\]
The positive sum has tail at most $1/(2N^2)$ after $N$ terms. Hence
$C_s$ can be enclosed using a finite sum and this explicit remainder.
Together with the endpoint moments and
\eqref{pf:polynomial-fourier-contract}, these bounds give the
archimedean inputs to Proposition~\ref{fw:transfer}.
Section~\ref{sec:certificate} supplies the finite-place inputs and
evaluates the resulting margin.

We also use the following digamma enclosure. For $x>0$, integers $L\ge0$
and $m\ge1$, and $z=x+L$,
\[
 \left|\psi(x)-\left(\log z-\frac1{2z}
 -\sum_{k=1}^{m-1}\frac{B_{2k}}{2kz^{2k}}
 -\sum_{j=0}^{L-1}\frac1{x+j}\right)\right|
 \le\frac{|B_{2m}|}{2m z^{2m}},
\]
where $B_{2k}$ are rational Bernoulli numbers. To verify the remainder,
expand $(z^2+t^2)^{-1}$ as a finite geometric sum in Binet's identity
\cite[Equation~(5.9.15)]{DLMF},
\[
 \psi(z)=\log z-\frac1{2z}
 -2\int_0^\infty\frac{t}{(t^2+z^2)(e^{2\pi t}-1)}\,dt.
\]
The residual has a fixed sign and is bounded by the next integrated term,
$|B_{2m}|/(2m z^{2m})$. Apply the digamma recurrence $L$ times.
We evaluate the displayed expression with exact rational arguments and
Bernoulli numbers, outward arithmetic, and the explicit remainder.
The rational witness values are used without rounding the arguments.

\section{The Rational Witness}
\label{sec:certificate}

We apply Proposition~\ref{fw:transfer} with the following rational
parameters and weights. Every parameter and weight written as a finite
decimal is exact. The displayed inequalities are rounded outwards from
more precise enclosures.

\subsection{Parameters and Local Weights}
Set
\[
 \delta=\frac{83647}{2000000},\quad p=\frac2{1+\delta},\quad
 C=\frac{42161819}{10^9},\quad
 \theta_* =\frac{4095}{8192},\quad \varepsilon=10^{-9}.
\]
The real-place Gaussian is $f_C(z)^p=e^{-a_C|z|^2}$ with $a_C=2p\delta$.
At a complex pair, use the polynomial profile and coefficient matrix in
\eqref{pf:bernstein-witness}, with $s=1.14600585$ and
$a=0.0000128633241758$. For the Fourier estimate, take $\rho=10^{-3}$.
For the mass integral, use an $8\times8$ rational subdivision and
binomial degree $N=18$ (indices $0,\ldots,18$). For the overlap, retain
orders zero and one with the signed remainder proved above.

At each selected rational prime $r$, let $Q=r^{f_r}$ and use identical
six-shell profiles on the two local coordinates, with weights
$w_{r,0},\ldots,w_{r,5}$. All $w_{r,0}=1$. The absolute local indices
and additive periods are
\[
\begin{array}{c|rrrrrrrrrrr}
r&2&3&5&7&11&13&17&19&23&29&31\\\hline
e_r&8&2&2&2&2&2&1&1&1&1&1\\
f_r&4&2&2&4&4&4&4&4&4&4&4\\
k_r&7&9&6&2&2&1&1&1&1&1&1
\end{array}
\]
The weights in Tables~\ref{tab:weights123} and~\ref{tab:weights45} are
strictly positive and nonincreasing. Rational comparisons verify these
properties and positivity of the Bernstein coefficients. The finite
local factor is
\begin{equation}\label{cert:finite-factor}
 \log\mathcal F_{\delta,r}
 =-\delta k_r\log Q+
 \log\big((k_r+1)B_r^2+2B_rR_r\big)-2(1+\delta)\log A_r,
\end{equation}
where $A_r=\sum_i m_iw_{r,i}^p$, $B_r=\sum_i m_iw_{r,i}^2$,
$R_r=\sum_{i,j}D_{ij}w_{r,i}w_{r,j}$, and the rational shell masses
$m_i$ and matrix $D$ are those of Section~\ref{fw:section}.
Expanding the shell weights into a nonnegative sum of ball indicators
gives a second calculation of \eqref{cert:finite-factor}. The overlap
numerators agree exactly as rational numbers.

\begin{table}[htbp]
\centering\scriptsize
\caption{Exact weights for shells 1, 2 and 3. All $w_{r,0}=1$.}
\label{tab:weights123}
\begin{tabular}{rrrr}
\toprule
$r$ & $w_{r,1}$ & $w_{r,2}$ & $w_{r,3}$\\
\midrule
2 & $5.2903102784782\cdot10^{-2}$ & $2.3157958402044\cdot10^{-3}$ & $1.0125841840494\cdot10^{-4}$\\
3 & $9.8866787345973\cdot10^{-2}$ & $8.1985230138641\cdot10^{-3}$ & $6.7798519183202\cdot10^{-4}$\\
5 & $3.1389392203562\cdot10^{-2}$ & $8.3331551426164\cdot10^{-4}$ & $2.2242307460119\cdot10^{-5}$\\
7 & $2.1902148915124\cdot10^{-4}$ & $3.888831805136\cdot10^{-8}$ & $8.3349312286089\cdot10^{-14}$\\
11 & $2.2066794391725\cdot10^{-5}$ & $5.2893244151477\cdot10^{-10}$ & $3.2351541879579\cdot10^{-16}$\\
13 & $1.9753217475271\cdot10^{-5}$ & $4.7825791524842\cdot10^{-11}$ & $1.154930945714\cdot10^{-16}$\\
17 & $5.1855578363751\cdot10^{-6}$ & $4.661704085405\cdot10^{-12}$ & $4.1907708339914\cdot10^{-18}$\\
19 & $3.0150737786821\cdot10^{-6}$ & $1.7755448176562\cdot10^{-12}$ & $1.0455994096796\cdot10^{-18}$\\
23 & $1.2033271761207\cdot10^{-6}$ & $3.17236260693\cdot10^{-13}$ & $8.3633817213758\cdot10^{-20}$\\
29 & $4.0183950841977\cdot10^{-7}$ & $4.2255710985501\cdot10^{-14}$ & $4.4434284649605\cdot10^{-21}$\\
31 & $2.9331492202408\cdot10^{-7}$ & $2.5364272697667\cdot10^{-14}$ & $4.0395513941639\cdot10^{-21}$\\
\bottomrule
\end{tabular}
\end{table}

\begin{table}[htbp]
\centering\scriptsize
\caption{Exact weights for shells 4 and 5. All $w_{r,0}=1$.}
\label{tab:weights45}
\begin{tabular}{rrr}
\toprule
$r$ & $w_{r,4}$ & $w_{r,5}$\\
\midrule
2 & $4.4907846711318\cdot10^{-6}$ & $2.0159593157185\cdot10^{-7}$\\
3 & $5.6589534354296\cdot10^{-5}$ & $4.7456521092252\cdot10^{-6}$\\
5 & $6.0504358878881\cdot10^{-7}$ & $1.6737711273851\cdot10^{-8}$\\
7 & $1.6239748423571\cdot10^{-18}$ & $3.1643053340993\cdot10^{-23}$\\
11 & $1.2374208229268\cdot10^{-21}$ & $4.7330769061366\cdot10^{-27}$\\
13 & $2.7890532331439\cdot10^{-22}$ & $6.7353099910173\cdot10^{-28}$\\
17 & $3.7674120582188\cdot10^{-24}$ & $3.3868217038473\cdot10^{-30}$\\
19 & $6.1574234266266\cdot10^{-25}$ & $3.6260409965598\cdot10^{-31}$\\
23 & $2.2048599887241\cdot10^{-26}$ & $5.812729505639\cdot10^{-33}$\\
29 & $4.6725178828482\cdot10^{-28}$ & $4.9134184420207\cdot10^{-35}$\\
31 & $6.4334537606741\cdot10^{-28}$ & $1.0246020721641\cdot10^{-34}$\\
\bottomrule
\end{tabular}
\end{table}


\subsection{Certified Enclosures}
Let $q=sp-1$, and let $T,U$ be independent with density
$qt^{q-1}\,dt$ on $(0,1)$. The mass and overlap formulas of
Section~\ref{pf:section} give
\begin{align}
38.9687816563247549
 &<\mathbb E[P(T,U)^p]<38.9687816563247561,
 \label{cert:mass}\\
348.4186894704059951
 &<Z_*<348.4186894704059952.
 \label{cert:overlap}
\end{align}
Here $Z_*$ is the lower expression for the overlap used in
\eqref{pf:collected-functional}. The upper endpoint in
\eqref{cert:overlap} bounds this expression and gives no upper bound
for the overlap integral. Write $J_D^*$ for the right side of
\eqref{pf:collected-functional} with the certified mass allowance. Thus
$J_D\ge J_D^*$. The finite terms satisfy
\[
 1.0335669225039939710
 <\sum_r\frac{\log\mathcal F_{\delta,r}}{e_rf_r}
 <1.0335669225039939712.
\]
The rational polynomial tube bound satisfies $q_0<0.026<1$.
The Fourier and period calculation yields
\[
 \log K_D<0.04916819,\qquad \sigma_D>1.7518755,
 \qquad \mu>26497.47.
\]
Here $\sigma_D=2\pi\rho/\sqrt a$, and $K_D$ is the Fourier prefactor
in \eqref{pf:polynomial-fourier-contract}. Together with the Gaussian
bound, these inequalities imply \eqref{fw:fourier-condition}.
The endpoint moments are finite and the common energy window is
compact by Sections~\ref{geo:section}, \ref{fw:section}
and~\ref{pf:section}.

Define the conservative margin
\begin{align*}
 M_*(\theta)={}&\sum_r\frac{\log\mathcal F_{\delta,r}}{e_rf_r}
 -(1/2-\delta)\ell-C+(1-\delta)\log2+(1-\theta)\log\pi\\
 &+(1-2\theta)J_C+\theta J_D^*.
\end{align*}
Since $\theta\ge0$, this is a lower bound on
\eqref{fw:margin}. Its signature slope is
$-\log\pi-2J_C+J_D^*>0.6492390240$. It suffices to evaluate the margin
at the smallest permitted signature. Directed interval arithmetic gives
\begin{align}
 M_*(\theta_*)&>0.00000533604359816794,
 \label{cert:margin}\\
 M_*(\theta_*)-4\varepsilon
 &>0.00000533204359816794>0.
 \label{cert:final-margin}
\end{align}
These evaluations use the exact inputs and full interval expressions.
The abbreviated decimal bounds displayed here are not substituted into
later calculations.

\begin{proof}[Proof of Theorem~\ref{thm:main}]
Theorem~\ref{tw:field-family} gives finite Galois fields $K_j$ of unbounded
degree with the prescribed root discriminant and local indices. Their
actual complex conjugations define quadratic extensions $K_j/F_j$
unramified at every finite place, with all selected primes split and
$\theta_j\ge\theta_*$. Theorem~\ref{an:ceiling} gives
$d_j^{-1}\log L_{F_j}(1)<C$ for all sufficiently large $j$.

The profiles above satisfy the positivity, period, moment and Fourier
conditions of Proposition~\ref{fw:transfer}. Their margin is at least
$M_*(\theta_*)$, and \eqref{cert:final-margin} is positive. The
proposition gives finite planar sets whose unordered edge ratio and
cardinalities tend to infinity, proving Theorem~\ref{thm:main}.
\end{proof}

\subsection{The Supplementary Certificate}
The rational witness is \path{certificates/witness.json}. The source and arithmetic
data are distributed as the three parts
\path{certificates/data.tar.gz.part00} through \path{part02}.
The replay script verifies and joins the parts. The part and
member SHA-256 index is \path{certificates/data-manifest.json}.
The archive includes the
coefficient generators and complete retained moment tables, actual-field
arithmetic receipts, and the prime census through $10^{12}$. Its internal
directory names record computational provenance. The proofs are given
in this manuscript.

Finite-field linear algebra verifies the quadratic forms, ranks and
representation counts. Exact Euler recurrences enumerate signed and
absolute coefficients. Integer bin moments, the explicit kernel
remainders, and outward arithmetic enclose the Hecke values. A prime
sieve supplies the finite Frobenius correction. Rational polynomial
operations and the mass and overlap remainders of
Section~\ref{pf:section} enclose the geometric score. The preceding
sections justify these calculations. The verifier evaluates the fixed
witness and contains no optimizer.

The replay works on a private copy of the supplied inputs, which it
leaves unchanged, and has three levels. The basic level checks the
retained input identities and recomputes the final geometric witness.
The analytic level also evaluates every Hecke factor at $12001/12000$,
checks the independent prime-census prefix and the full finite
correction, and reassembles the Euler ceiling. The moment level also
regenerates the 64 complete sector files of the five newest sparse
producers and compares every integer bin; these contain 103012 bins and
5515960 integer moments. Earlier moment tables and actual-field
arithmetic receipts remain supplied finite data at every level, and no
level repeats the full $10^{12}$ prime sieve, whose source is included.

The arithmetic software uses PARI/GP for exact field and Euler
computations, Arb for ball arithmetic, and MPFR and mpmath intervals for
directed enclosures \cite{PARI,Johansson2017,MPFR,Mpmath}.
The package README gives the replay commands, dependency versions and
build instructions. A manifest binds the bytes of every source file and
the bibliography of this manuscript together with the certificate
archive, and its checker rejects any change to that set. These
checks establish source identity. The mathematical proofs and the
finite computations remain separate obligations.

\clearpage
\bibliographystyle{amsplain}
\bibliography{references}
\end{document}
